\documentclass[10pt]{article} % For LaTeX2e
\usepackage{tmlr}
% If accepted, instead use the following line for the camera-ready submission:
%\usepackage[accepted]{tmlr}
% To de-anonymize and remove mentions to TMLR (for example for posting to preprint servers), instead use the following:
%\usepackage[preprint]{tmlr}

% Optional math commands from https://github.com/goodfeli/dlbook_notation.
\input{math_commands.tex}

\usepackage{hyperref}
\usepackage{url}
\usepackage{graphicx}
\usepackage{booktabs}
\usepackage{multirow}
\usepackage{amsmath,amssymb,amsthm}
\usepackage{xcolor}
\usepackage{array}
\usepackage{tabularx}

\newtheorem{proposition}{Proposition}
\newtheorem{corollary}{Corollary}
\newtheorem{remark}{Remark}
\newcommand{\nom}{nearest-of-$m$}
\newcommand{\Nom}{Nearest-of-$m$}
\newcommand{\imlewm}{IMLE-WM}
\newcommand{\ms}[2]{$#1 \pm #2$}
\newcommand{\best}[1]{\textbf{#1}}
\newcommand{\dmin}{d_{\min}}
\newcommand{\deff}{d_{\mathrm{eff}}}
\newcommand{\Imk}{I(m,K)}
\newcommand{\head}[1]{\multicolumn{1}{c}{#1}}
\newcommand{\IS}{\mathrm{ES}}
\graphicspath{{figures/}}

\title{How Many Candidates Does a One-Step World Model Need?\\A Random-Codebook Account of \Nom{} Training on Frozen DINO Tokens}

% Authors must not appear in the submitted version. They should be hidden
% as long as the tmlr package is used without the [accepted] or [preprint] options.
% Non-anonymous submissions will be rejected without review.

\author{\name Anonymous authors \email \\
      \addr Paper under double-blind review}

\def\month{MM}  % Insert correct month for camera-ready version
\def\year{YYYY} % Insert correct year for camera-ready version
\def\openreview{\url{https://openreview.net/forum?id=XXXX}} % Insert correct link to OpenReview for camera-ready version

\begin{document}

\maketitle

\begin{abstract}
Deterministic world models that predict future vision-foundation-model features average over multimodal futures; stochastic alternatives recover the modes but usually pay tens of network evaluations per step. A minimal one-step alternative trains a noise-conditioned generator with a \nom{} objective (conditional implicit maximum likelihood estimation, IMLE): draw $m$ candidate futures and pull only the one closest to the observed outcome. Recent work conjectures that a small candidate budget suffices for feature-space world models and, conversely, that a strictly proper energy score should be preferred over such hard-minimum objectives. We test both claims with action-conditioned one-step world models on frozen DINOv2/v3 patch tokens, on benchmarks whose futures can be enumerated. Under mode-like uncertainty, \nom{} with $m{=}16$ covers the enumerated futures better than 50-step flow matching, one-step flow generators (iMF, Shortcut) and an energy-score generator at matched parameters, with about $30\times$ lower latency, and its autoregressive rollouts stay on the token manifold where flow rollouts drift regardless of the number of integration steps. Our main results concern the budget. The \nom{} objective is the expected distortion of a random codebook of $m$ draws. This view predicts, and we confirm with predictions recorded before the runs, three things: (i) in contexts that never repeat, the generator memorizes the realized training outcome, and the strength of this memorization is governed by the codebook's miss \emph{distortion} $I(m)$ rather than by the miss probability (seven of nine predictions across $K{=}2$--$16$ modes fall in their intervals), whereas an energy-score generator memorizes regardless of $m$ because its incentive is $m$-independent; (ii) mode probabilities flatten with $m$, matching a known optimum within $0.01$; (iii) in high-dimensional continuous outcome spaces the samples under-disperse, which bounds the approach on real robot video. Per-step matching on self-rollouts reduces long-horizon error by a third on a stochastic Push-T variant, and an ablation separates its two effects: any self-generated context keeps rollouts on the manifold, but only feeding back the candidate nearest the truth improves long-horizon coverage. Planning success does not improve despite accurate branch-risk estimates, a negative result we report in full.
\end{abstract}

\section{Introduction}
\label{sec:intro}

World models that predict the future in the feature space of a frozen vision foundation model are cheap to train and plan with \citep{zhou2025dinowm,baldassarre2025dinoworld,assran2025vjepa2}. Their weakness is well known: a deterministic predictor trained with a squared loss returns the conditional mean, and when the future has several modes the mean lies between them, blurring the prediction and compounding over a rollout \citep{porcher2026flowwm,baldassarre2025dinoworld}. Stochastic feature-space world models fix this with diffusion or flow matching \citep{porcher2026flowwm,boduljak2025vfmf,walker2025generalist}, at the price of tens of network evaluations (NFE) per predicted step, which matters when a planner must imagine thousands of futures per decision. A growing line of work therefore builds one-step stochastic world models: drifting objectives \citep{lu2026driftworld}, mean-flow heads \citep{agibot2026geact}, terminal maps \citep{lee2026beckmann}, priming \citep{park2026priming}, or streaming flows with one velocity evaluation per step \citep{hsu2026vsf}.

The simplest one-step objective is older than all of these. Draw $m$ candidate futures from a noise-conditioned generator, find the one closest to the observed outcome, and pull only that one towards it. This is conditional implicit maximum likelihood estimation \citep[IMLE;][]{li2018imle,li2019diverse}, the ``min-of-$K$'' objective of explorative modeling \citep[XM;][]{gladstone2026xm}, and a sampled relative of winner-takes-all and multiple-choice learning \citep{guzman2012mcl,rupprecht2017mhp}. Every data point is explained by at least one sample, so mode collapse is discouraged by the objective itself, and inference is one forward pass. IMLE world models already outperform Gaussian ones on proprioceptive state \citep{aghabozorgi2026wimle}, and \citet{gladstone2026xm} conjecture that feature-space world models have few modes, so a small candidate budget should suffice. In the opposite direction, \citet{makabe2026esp} argue that a hard minimum over a finite number of candidates is not a strictly proper scoring rule and show that a one-step generator trained with the energy score covers a two-branch toy distribution better than the min-of-$K$ objective.

This paper asks a quantitative version of both claims: \emph{how many candidates does a one-step \nom{} world model need, when is no budget enough, and is a strictly proper score the better objective?} We answer with a conditional IMLE world model on frozen DINOv2/v3 patch tokens, trained and compared under matched parameters against deterministic regression, 50- and 2-step flow matching in the FlowWM recipe, two one-step flow generators (improved MeanFlow and Shortcut models), and an energy-score generator, on benchmarks whose futures can be enumerated so that mode coverage is measured rather than inferred.

One lens organizes the results. The \nom{} loss is, in expectation, the distortion incurred when the outcome is quantized with a \emph{random codebook} of $m$ independent draws from the generator. This correspondence is a sampled version of the known equivalence between winner-takes-all learning and conditional quantization \citep{rupprecht2017mhp,letzelter2024wta,perera2024amcl}, and its continuous limit is the Gaussian rate--distortion test channel \citep{shannon1959coding,cover2006elements}. We make no claim to new mathematics; the contribution is to show that this one correspondence predicts, quantitatively and before the experiments were run, where the candidate budget matters and where it cannot help.

\paragraph{Contributions.}
\begin{enumerate}
\item \textbf{A controlled test of the small-budget conjecture.} Under mode-like uncertainty (bouncing dynamics with $K$ enumerable outcomes, a Push-T variant with a hidden center of gravity), \nom{} with $m{=}16$ covers enumerated futures better than 50-step flow matching, the one-step flow generators and the energy-score generator at matched parameters (coverage $0.93$ vs.\ $0.79$ for FM-50 and $0.36$ for the energy score on non-repeating contexts; F1 error $7.2$ vs.\ $12.3$\,px on Bouncing Shapes), with about $30\times$ lower latency than FM-50 (Section~\ref{sec:results1}). On real robot video it does not, and Section~\ref{sec:results3} explains why.
\item \textbf{Outcome memorization and its incentive.} In contexts that never repeat, a conditional generator can reach zero training loss by memorizing the realized outcome of each training context, keyed on its appearance, and ignoring its noise input. We measure this with the share of samples that fall on the realized mode of training contexts (chance $1/K$), show that it is governed by the random codebook's miss \emph{distortion} $I(m)$ rather than by the miss probability $(1-1/K)^m$ (RMSE $0.031$ vs.\ $0.068$ over $K\in\{2,4,8,16\}$), and confirm seven of nine predictions recorded before the $K{=}2,16$ runs. The same accounting explains why the strictly proper energy score memorizes in this regime while \nom{} with $m{=}16$ does not: its incentive is half the expected outcome spread, independent of $m$ (Section~\ref{sec:results2}).
\item \textbf{Where no budget suffices.} The codebook view also predicts a flattening of mode probabilities with $m$, which we confirm within $0.01$ of the optimum derived by \citet{bennett2026freedom}, and an under-dispersion of $e_m/\sqrt{d}$ in $d$-dimensional continuous outcome spaces, which is the form of the Gaussian rate--distortion test channel. On BAIR robot pushing the samples spread along only one to three context-specific directions and no budget up to $m{=}64$ closes the gap to multi-step flow matching on distribution metrics (Section~\ref{sec:results3}).
\item \textbf{Rollouts and planning.} Autoregressive \nom{} rollouts stay as close to the token manifold as deterministic regression, while flow rollouts drift by a factor of four at $h{=}32$ regardless of the number of integration steps. Per-step matching on self-generated contexts, in the lineage of DaD and closest-among-top-$K$ rollouts \citep{venkatraman2015dad,zhang2025catk}, reduces long-horizon recall error by a third on stochastic Push-T; an ablation of the fed-back state separates its manifold benefit from its coverage benefit. Planning with a cross-entropy-method controller is $50\times$ faster than with FM-50 at equal success, but success does not improve even on a task where the model's branch-risk estimates are accurate (Section~\ref{sec:results4}).
\end{enumerate}

We also report what did not work: trajectory-level matching, multi-stage supervision in token space, variance-penalized planning and a counterfactual action-contrast term (Appendix~\ref{app:negative}). Benchmarks are synthetic or simulated except BAIR, models are small (8M and 21M parameters), and our flow-matching baselines are not at the scale of \citet{porcher2026flowwm}; Section~\ref{sec:discussion} states these limits.

\section{Related Work}
\label{sec:related}

\paragraph{Min-of-$m$ objectives.}
IMLE \citep{li2018imle} matches each data point to its nearest generated sample from a shared pool; conditional IMLE uses a per-condition pool and was found to ignore its noise input when each condition has a single target \citep{li2019diverse,li2020multimodal}. Later variants reject samples that are already close \citep[RS-IMLE;][]{vashist2024rsimle}, anneal a per-sample radius \citep{aghabozorgi2023adaptive}, split the noise hierarchically \citep{peng2022chimle}, or add multi-stage supervision and a robust loss \citep[ROMS-IMLE;][]{vashist2026roms}. IMLE has been used as a world model on state vectors \citep{aghabozorgi2026wimle}, as a trajectory generator for model-predictive control \citep{lee2026imlempc} and as a one-step action head \citep{rana2025imlepolicy,hosseinkhani2026imlevla}. Explorative modeling \citep{gladstone2026xm} studies the same min-of-$K$ objective as a pretraining axis, includes a one-step world model on Maze2D, and shows on video that increasing $K$ reduces overfitting; it explains this by expressivity, whereas we quantify the incentive. The conditional min-of-$N$ idea goes back to point-set generation \citep{fan2017pointset} and best-of-many sampling for sequence prediction \citep{bhattacharyya2018bestofmany}; multi-hypothesis prediction of future features was done with hard assignment by \citet{vondrick2016anticipating}. Winner-takes-all and multiple-choice learning \citep{guzman2012mcl,lee2016smcl,rupprecht2017mhp,makansi2019ewta,seo2020tmcl} learn $K$ heads instead of sampling $m$ candidates; their optima are conditional quantizers \citep{rupprecht2017mhp,letzelter2024wta,delattre2025cdq,cortes2025timemcl}, their training follows a rate--distortion curve \citep{perera2024amcl}, and their hard assignment distorts mode weights \citep{wu2026wtagmm}. \citet{thiede2019variety} show that an i.i.d.\ min-of-$N$ loss over-disperses a low-dimensional density, the regime opposite to the one we find on video. \citet{bennett2026freedom} derives the optimal mode weights of a min-of-$K$ objective; we confirm that result in a learned world model. \citet{makabe2026esp} compare the objective against the energy score \citep{gneiting2007strictly}; Section~\ref{sec:es} tests that comparison in our setting. Branch-JEPA \citep{song2026branchjepa} predicts a finite set of latent successors and compares specialization with energy-score training.

\paragraph{Stochastic and one-step world models in feature space.}
DINO-WM, DINO-world and V-JEPA~2 \citep{zhou2025dinowm,baldassarre2025dinoworld,assran2025vjepa2} predict frozen DINO features deterministically. FlowWM \citep{porcher2026flowwm} runs flow matching directly in DINOv3 feature space with a wide velocity head and a shifted time schedule, introduces the enumerable Bouncing Shapes benchmark that we adopt, and reports 50-step sampling; VFMF \citep{boduljak2025vfmf} and \citet{walker2025generalist} use diffusion on compressed or frozen features. One-step stochastic world models include DriftWorld \citep{lu2026driftworld}, the mean-flow planner of GE-Act 2.0 \citep{agibot2026geact}, Beckmann world models \citep{lee2026beckmann}, priming models \citep{park2026priming} and variational streaming flows \citep{hsu2026vsf}; none of them measures enumerable mode coverage, and \citet{lee2026beckmann} report no planning improvement on Push-T, consistent with our negative planning result. Our flow baselines follow iMF \citep{geng2025imf}, Shortcut models \citep{frans2025shortcut}, which also underlie Dreamer~4 \citep{hafner2025dreamer4}, and the FlowWM recipe. Generated latents leaving the manifold of clean encoder features is a known problem \citep{zheng2025rae,boduljak2025vfmf,karypidis2026latentforesight}; we measure it for rollouts of each objective.

\paragraph{Memorization in conditional generators.}
Conditional diffusion models memorize more when the condition is more informative, and random per-sample labels trigger memorization \citep{gu2025memorization,kim2026memorization}; data copying is detected by local over-representation of training points \citep{meehan2020nonparametric,somepalli2023diffusion,carlini2023extracting}. Conditional GANs ignoring their noise is a related collapse \citep{yang2019dsgan,mao2019msgan}. \citet{li2018imle} already note that with unlimited capacity any sample-based objective memorizes. Our measurement differs in that it labels outcomes by mode and compares training and held-out contexts, which isolates memorization of the realized outcome from a failure to generalize.

\paragraph{Rollout training and exposure bias.}
Training on self-generated context is the classic remedy for compounding error \citep{ross2011dagger,venkatraman2015dad,bengio2015scheduled,lamb2016professor,hafner2019planet}; recent video and latent world models use noise injection \citep{chen2024diffusionforcing}, self-generated prefixes with a teacher \citep{huang2025selfforcing} or bootstrapped latents \citep{gao2026sword}. Closest-among-top-$K$ fine-tuning \citep{zhang2025catk} rolls out a tokenized policy on the candidate nearest the ground truth. Our per-step matching is the continuous-feature analogue: in \nom{} training the nearest candidate is already the argmin of the loss. \citet{terver2025drives} find that multi-step rollout losses help planning only up to short horizons, and \citet{cheung2026fidelity} that one-step fidelity predicts planning poorly; our planning results agree.

\section{\Nom{} World Models on Frozen DINO Tokens}
\label{sec:method}

\subsection{Setting and notation}
A frozen encoder $E$ maps each observation $o_t$ to patch tokens $s_t = E(o_t) \in \R^{N \times D}$ with $N = 256$ tokens on a $16\times16$ grid and $D = 384$ channels. For the bounce, push, Push-T and BAIR benchmarks $E$ is DINOv2 ViT-S/14 at $224$\,px \citep{oquab2024dinov2}; for Bouncing Shapes it is the DINOv3 ViT-S/16 feature of FlowWM \citep{simeoni2025dinov3,porcher2026flowwm}, the average of the normalized hidden states of layers 3, 6, 9 and 12 at $256$\,px. Tokens are precomputed once, stored in half precision and standardized per token. The world model takes a context of $k = 3$ token maps $c_t = (s_{t-2}, s_{t-1}, s_t)$, an action $a_t$ and a noise vector $z \in \R^{256}$, and predicts the next token map as a residual,
\begin{equation}
\hat{s}_{t+1} = s_t + g_\theta(c_t, a_t, z).
\end{equation}
Rollouts feed $\hat{s}_{t+1}$ back into the context; every step is one network evaluation.

\subsection{Generator}
$g_\theta$ is a ConvNeXt stack \citep{liu2022convnext} on the $16\times16$ token grid. The context maps are concatenated along channels and projected by a $1\times1$ convolution; the action passes through a two-layer MLP and modulates every block by FiLM; the noise $z$ passes through a four-layer MLP and modulates blocks by adaptive layer normalization, as in the mapping network of ROMS-IMLE \citep{vashist2026roms}. The output head is a per-token MLP of width $1024$, the ``wide head'' that FlowWM found necessary for velocity prediction in feature space; we attach the same head to every method so that parameter counts match within $\pm4\%$. We use two sizes: width 256 with 8 blocks ($8.0$--$8.2$M parameters across methods) for the synthetic benchmarks, and width 384 with 12 blocks ($20.6$M) for Push-T and BAIR. The noise enters only the last three (respectively four) blocks, so the context-processing trunk is evaluated once per transition and only the tail is evaluated per candidate; this \emph{late fusion} halves training time at equal coverage (Appendix~\ref{app:impl}).

\subsection{Objective: \nom{} matching with a robust loss}
\label{sec:objective}
For each transition $(c_t, a_t, s_{t+1})$ in a batch we draw $m$ noise vectors, generate $m$ candidates $\hat{s}^{(j)} = s_t + g_\theta(c_t, a_t, z_j)$ without gradient, and select the one nearest to the realized outcome in token space,
\begin{equation}
j^\star = \argmin_{j \in [m]} \big\| s_{t+1} - \hat{s}^{(j)} \big\|^2 ,
\label{eq:match}
\end{equation}
where $\|\cdot\|$ is the mean squared error over all tokens and channels. The selected candidate is regenerated with gradient and trained with a Geman--McClure loss \citep{geman1987statistical} on its per-sample root-mean-square distance $d$,
\begin{equation}
\mathcal{L} = \rho(d),\qquad \rho(d) = \frac{2u^2}{u^2 + 4},\quad u = \frac{d}{2c},
\label{eq:gm}
\end{equation}
with the scale $c$ tracked as the running median of the matching distance. The matching is recomputed for every batch (a per-transition pool, as in WIMLE; a shared pool as in image IMLE would be meaningless for conditional prediction), runs in bfloat16 under no-gradient and, importantly, in its own autocast region: autocast caches weight casts, and casts made under no-gradient and reused by the gradient pass silently cut the gradient to the weights. The robust loss reproduces the ROMS-IMLE observation in our pilot (precision $0.92$ vs.\ $0.86$ for $\ell_2$ with the pilot recipe) but matters little in the final recipe (Appendix~\ref{app:ablation}); the candidate budget $m$ is the variable this paper is about. We call the model \imlewm{} and write $m$ for the budget throughout.

\paragraph{Memorization monitor.}
Every 2{,}000 steps we draw 16 samples for 200 training contexts and 200 held-out contexts whose next transition is stochastic, decode them with a position readout, and record the share that falls on the branch that was actually realized in the data. For a held-out context this share is $1/K$ by construction; for a training context it exceeds $1/K$ only if the model has learned which of the $K$ branches that particular context led to. This \emph{memorization index} is the measurement behind Section~\ref{sec:results2} and doubles as a cheap early warning during training.

\subsection{Per-step matching on self-rollouts}
\label{sec:seqmatch}
A world model is trained on real contexts and used on its own predictions. The natural IMLE remedy, matching whole $H$-step candidate trajectories, fails: the number of trajectory modes grows like $K^H$, far beyond any budget $m$, so the selected candidate is the one that is on average close to all branches, and averaging returns (Appendix~\ref{app:negative}). We instead match one step at a time along the model's own rollout. Starting from a real context, at step $h = 1,\dots,H$ we draw $m$ fresh noise vectors, select the candidate nearest to the true frame $s_{t+h}$ \emph{given the self-predicted context}, train that step with Equation~\ref{eq:gm}, and feed the selected prediction (detached) back as the next context. With $H = 4$ this costs $mH$ forward passes without gradient and $H$ with gradient per transition. The fed-back state has two natural alternatives, a random candidate (a free-running rollout) and the true frame (teacher forcing with the same number of matched updates); Section~\ref{sec:seq} uses them to separate what the procedure does. The selection of the candidate nearest the truth is the same device as closest-among-top-$K$ rollouts \citep{zhang2025catk}; in \nom{} training it is the argmin of the loss itself and needs no extra model.

\subsection{Inference and planning}
One predicted step is one forward pass; $m_p$ futures for a state are obtained by batching $m_p$ noise vectors. For planning we use cross-entropy-method model-predictive control \citep[CEM-MPC;][]{rubinstein1999cem,zhou2025dinowm}: each candidate action sequence is imagined with $m_p$ noise draws and scored by the mean, over draws, of a distance between the imagined terminal state and the goal. We do not penalize the spread of the imagined futures; a mean-plus-variance cost lowered the success rate even with the true simulator as the world model (Appendix~\ref{app:negative}).

\subsection{Baselines at matched parameters}
\label{sec:baselines}
All baselines share the backbone, the wide head, the context, the residual parameterization and the training budget unless stated. \emph{Regression} minimizes the squared error of a single prediction. \emph{FM-50} is flow matching in the FlowWM recipe \citep{lipman2023flow,porcher2026flowwm}: the velocity head predicts the flow between Gaussian noise and the normalized residual, with the time shift $t' = \alpha t/(1 + (\alpha-1)t)$, $\alpha = 13$, applied to both training and the 50-step Euler sampler; \emph{FM-2} is the same model sampled with 2 Euler steps. \emph{iMF} is improved MeanFlow \citep{geng2025meanflow,geng2025imf} with its dual average- and instantaneous-velocity heads, ported from the official implementation and trained for three times as many steps because it was still improving at the common budget. \emph{Shortcut} is a shortcut model \citep{frans2025shortcut} with a quarter of each batch devoted to the self-consistency target, trained for twice as many steps so that its training time matches \imlewm{} with $m{=}16$, and sampled in one step (two steps changed its metrics by less than $0.01$). \emph{Energy score} trains the \imlewm{} generator unchanged with the energy score of \citet{makabe2026esp}: $m$ samples per transition, the unbiased estimator $\frac{1}{m}\sum_j \|Y_j - x\|^\beta - \frac{1}{2}\binom{m}{2}^{-1}\sum_{i<j}\|Y_i - Y_j\|^\beta$ with $\beta = 1$ and the whole-map root-mean-square distance, with $m \in \{8, 16\}$ and a per-token variant that averages a per-token score over the grid. Parameter counts, schedules and the porting notes for iMF are in Appendix~\ref{app:impl}.

\section{A Random-Codebook Account of \Nom{} Training}
\label{sec:theory}

Fix a context $c$, write $p(\cdot\,|\,c)$ for the true outcome distribution and $q(\cdot\,|\,c)$ for the generator's. Because the matching in Equation~\ref{eq:match} draws the candidates independently from $q$, the population objective behind \nom{} training is
\begin{equation}
L(q) = \E_{x \sim p}\, \E_{Y_1,\dots,Y_m \sim q}\, \min_{j \in [m]} \| Y_j - x \|^2 ,
\label{eq:codebook}
\end{equation}
the expected distortion of quantizing $x$ with a random codebook of $m$ i.i.d.\ draws from $q$. Winner-takes-all learning with $K$ fixed hypotheses is the deterministic version of this problem, and its minimizers are centroidal Voronoi quantizers \citep{rupprecht2017mhp,letzelter2024wta}; for a Gaussian source the random-codebook version is the rate--distortion setting of \citet{shannon1959coding} with rate $\ln m / d$. The derivations below are elementary consequences of Equation~\ref{eq:codebook}; their value is that they make predictions that can be checked, and in Sections~\ref{sec:results2} and~\ref{sec:results3} they were checked against predictions written down before the corresponding runs. Proofs are in Appendix~\ref{app:proofs}.

\subsection{Outcome memorization in non-repeating contexts}
\label{sec:theory_memo}
Let the outcome of context $c$ take one of $K$ modes $\mu_1,\dots,\mu_K$ with probabilities $p_1,\dots,p_K$. In training data whose contexts never repeat, each training context $c_i$ is seen with exactly one realized mode $\mu_{k_i}$. Two solutions compete for the empirical loss. The \emph{memorizing solution} $q(\cdot\,|\,c_i) = \delta_{\mu_{k_i}}$ puts all candidates on the realized mode and attains zero loss, at the cost of storing $k_i$ for every training context, which requires a feature of $c_i$ that identifies it. The \emph{generalizing solution} $q(\cdot\,|\,c) = p(\cdot\,|\,c)$ is the right answer on new contexts and attains
\begin{equation}
I(m) = \E_{k \sim p}\, \E_{S_1,\dots,S_m \sim p}\, \min_{j} \| \mu_{S_j} - \mu_k \|^2 ,
\label{eq:I}
\end{equation}
the expected distortion of a random codebook of $m$ mode draws.

\begin{proposition}[memorization incentive]
\label{prop:memo}
With non-repeating contexts, the empirical \nom{} loss prefers the memorizing solution to the generalizing one by exactly $I(m)$. If the modes are separated by at least $\delta$ and at most $\Delta$, then $(1 - 1/K)^m \delta^2 \le I(m) \le (1-1/K)^m \Delta^2$ for equiprobable modes, so the incentive vanishes exponentially once $m$ exceeds $K$. If instead the modes densely fill a curve of length $\ell$ with $K \gg m$, then $I(m) = \Theta(\ell^2/m^2)$, independent of $K$. For squared-error regression the same comparison gives $\mathrm{tr}\,\mathrm{Cov}_p$, which no budget reduces.
\end{proposition}

Two readings matter for what follows. First, the incentive is a \emph{distortion}, not a probability: it weighs a missed mode by how far away the nearest drawn mode is. Second, whether $K$ enters at all depends on the geometry. In our bounce benchmark the $K$ modes are equiangular directions on a semicircle, so for $m \lesssim K$ the distortion is nearly the same for every $K$ and only the position on the $I$ axis matters; the competing summary ``$m \approx 2K$'' that we had drawn from $K \in \{4, 8\}$ alone is a coincidence of those two points, and Section~\ref{sec:collapse} rejects it. We compute $I(m,K)$ for the benchmark geometry exactly (Appendix~\ref{app:imk}).

\subsection{High-dimensional continuous outcomes: shrinkage}
\label{sec:theory_shrink}
When the outcome is continuous and high-dimensional no codebook of $m$ draws can cover it, and the optimal $q$ shrinks.

\begin{proposition}[isotropic shrinkage]
\label{prop:shrink}
Let $p = \mathcal{N}(\mu, \sigma^2 I_d)$ and let the generator be free only in its spread, $q = \mathcal{N}(\mu, s^2 I_d)$. For large $d$,
\begin{equation}
L(q) \approx d\sigma^2 + d s^2 - 2 s \sigma \sqrt{d}\, e_m,\qquad e_m = \E\max_{j \le m} \xi_j,\ \xi_j \sim \mathcal{N}(0,1),
\end{equation}
so the optimal spread is $s^\star / \sigma = e_m / \sqrt{d}$ with $e_4 = 1.03$, $e_{16} = 1.77$, $e_{64} = 2.34$ and $e_m \approx \sqrt{2 \ln m}$ for large $m$.
\end{proposition}

This is the first-order form of the reproduction variance $\sigma^2(1 - m^{-2/d})$ of the Gaussian rate--distortion test channel at rate $\ln m / d$ nats per dimension \citep{cover2006elements}; the correspondence is exact because IMLE draws its codebook i.i.d.\ from the reproduction distribution. The practical content is a dichotomy in $\ln m$ against the effective dimension $\deff$ of the conditional outcome. If the outcome has a few discrete modes, $\ln m > \ln K$ lets the codebook cover them and there is no shrinkage; this is the regime of Proposition~\ref{prop:memo}, and it is also the regime in which an i.i.d.\ min-of-$N$ loss over-disperses a low-dimensional density \citep{thiede2019variety}. If $\deff \gg e_m^2$ the samples contract towards the conditional mean and the spread grows only like $\sqrt{\ln m}$. Section~\ref{sec:bair} shows that real robot video is in the second regime, and that the contraction is not isotropic.

\subsection{Flattening of mode probabilities}
\label{sec:theory_flatten}
Even when every mode is covered, the \nom{} optimum does not reproduce the mode probabilities. For equidistant modes the population optimum satisfies the stationarity condition $p_k (1 - q_k)^{m-1} = \mathrm{const}$ for all $k$ \citep[Theorem~7]{bennett2026freedom}; for $K = 2$ this gives $\big((1-q)/q\big)^{m-1} = (1-p)/p$. At $m = 1$ the solution collapses onto the most frequent mode (mode seeking, as for regression); at $K{=}2$, $m{=}2$ it is exactly calibrated; for large $m$ the differences $q_k - q_l \approx \ln(p_k/p_l)/(m-1)$ vanish, so rare modes are over-represented. For \emph{equiprobable} modes the optimum is exactly $q = p$ for every $m \ge 2$ \citep[Theorem~4]{bennett2026freedom}. The budget $m$ therefore trades coverage against calibration, and benchmarks with equiprobable branches cannot show it; Section~\ref{sec:flatten} uses a benchmark with $p = (0.8, 0.2)$.

\subsection{The energy score's incentive does not depend on $m$}
\label{sec:theory_es}
\citet{makabe2026esp} prefer the energy score $\IS_\beta(q, x) = \E_{Y \sim q}\|Y - x\|^\beta - \frac{1}{2}\E_{Y, Y' \sim q}\|Y - Y'\|^\beta$, $\beta \in (0, 2)$, over the hard minimum because the former is strictly proper. Propriety is a population property. In the finite-sample regime of Section~\ref{sec:theory_memo} it has a different consequence.

\begin{proposition}[memorization incentive of the energy score]
\label{prop:es}
(a) For fixed $x$, $\IS_\beta(q, x)$ is half the energy distance between $q$ and $\delta_x$, so it is non-negative and zero only at $q = \delta_x$ \citep{dawid2007geometry,gneiting2007strictly}. With non-repeating contexts the unique minimizer of the empirical energy-score risk is therefore the memorizing solution, for every number of samples $m \ge 2$. (b) The generalizing solution attains $\frac{1}{2}\E_{Y, Y' \sim p}\|Y - Y'\|^\beta$, the generalized entropy of the score, so the memorization incentive is half the expected outcome spread and does not depend on $m$; the $m$-sample estimator is unbiased, and $m$ only changes its variance. (c) For $\beta = 2$ the score reduces to squared-error regression of the mean. (d) Averaging a per-token energy score over the token grid inherits (a) and (b) token by token.
\end{proposition}

Dividing each incentive by the incentive of the corresponding fit-only objective gives a unit-free comparison: regression and the energy score sit at $1/2$ for every $m$, while \nom{} sits at $I(m)/I(1)$, which for $K{=}4$ equiprobable modes on the bounce geometry is $0.52, 0.21, 0.054, 0.005$ for $m = 2, 4, 8, 16$. The mathematics in (a)--(c) is classical and the unconditional case is already in \citet{li2018imle}; what is new is the comparison and its test in Section~\ref{sec:es}. A second difference concerns the direction of spread. With an unknown bias $a$ and an isotropic spread $s$ in $d$ dimensions, the energy score with $\beta < 2$ is minimized at $s = a$ with gain factor $2^{\beta/2 - 1}$ independent of $d$, a special case of Proposition~4.1 of \citet{debortoli2025distributional}, whereas Proposition~\ref{prop:shrink} gives $s \approx a\, e_m/\sqrt{d} \to 0$ for \nom{}. The energy score calibrates the \emph{size} of the spread but, in high dimensions, is insensitive to its \emph{direction}: the Euclidean energy distance asymptotically detects only differences in means and covariance traces \citep{chakraborty2021framework}, and the score is far more sensitive to errors in the mean than in the variance or dependence structure \citep{pinson2013discrimination,scheuerer2015variogram,marcotte2023regions}. Isotropic spread in token space leaves the data manifold; \nom{} under-disperses but stays on it. This part is an interpretation; Appendix~\ref{app:proofs} states its assumptions.

\subsection{Predictions tested in this paper}
\label{sec:predictions}
\begin{enumerate}
\item[P1] In non-repeating contexts, excess memorization $(\text{index} - 1/K)/(1 - 1/K)$ is a monotone function of $I(m,K)$ and collapses across $K$ when plotted against $I$, not against $(1-1/K)^m$ or $m/K$ (Prop.~\ref{prop:memo}; Section~\ref{sec:collapse}; predictions for $K{=}2,16$ recorded before the runs).
\item[P2] The energy score memorizes in non-repeating contexts at every $m$, and faster at larger $m$; it stops memorizing when contexts carry no identifying feature or are rarely revisited (Prop.~\ref{prop:es}; Section~\ref{sec:es}; predictions for the fixed-appearance and per-token runs recorded before the runs).
\item[P3] With $p = (0.8, 0.2)$ the share of the frequent mode among \nom{} samples falls with $m$ as the stationarity condition prescribes, and is calibrated for FM-50 (Section~\ref{sec:flatten}).
\item[P4] On real video the spread-to-error ratio of the samples grows with $m$ only like $e_m$ and saturates; raising $m$ to 64 does not reach the distribution metrics of 50-step flow matching (Prop.~\ref{prop:shrink}; Section~\ref{sec:bair}).
\item[P5] An energy-score generator spreads off the token manifold in directions that readouts do not see, so its advantage disappears under a noise-robust readout while \nom{} is unaffected (Section~\ref{sec:es}).
\end{enumerate}

\section{Experimental Setup}
\label{sec:setup}

\subsection{Benchmarks}
\label{sec:bench}
Table~\ref{tab:bench} summarizes the five settings; Appendix~\ref{app:bench} gives construction details.

\begin{table}[t]
\caption{Benchmarks. ``Contexts repeat'' means that the same context appears in the training data with several different outcomes, which removes the possibility of memorizing the realized outcome.}
\label{tab:bench}
\begin{center}
\small
\setlength{\tabcolsep}{4pt}
\resizebox{\linewidth}{!}{%
\begin{tabular}{llllll}
\toprule
Name & Source & Stochasticity & Contexts repeat & Actions & Model \\
\midrule
A1 Bouncing Shapes & FlowWM generator & normal bounce vs.\ reversal ($K{=}2$ per collision) & yes & none & 8M (joint: 27--29M) \\
A2-div / A2-fix bounce & ours & $K \in \{2,4,8,16\}$ directions at a wall; asymmetric $K{=}2$ & no / partly & none & 8M \\
Push (branch risk) & ours & $\pm20^\circ$ (2 modes), ice $\{-60^\circ,0,60^\circ\}$ (3), absorbing pit & no & 2-D push & 8M \\
B stochastic Push-T & DINO-WM data, re-simulated & hidden centre of gravity, 4 modes & no & 10-D (5 frames) & 21M \\
BAIR robot pushing & \citet{ebert2017bair} & continuous, high-dimensional & no & 4-D & 21M \\
\bottomrule
\end{tabular}}
\end{center}
\end{table}

\paragraph{A1: Bouncing Shapes.}
We use the public FlowWM generator unchanged \citep{porcher2026flowwm}: 32-frame clips at $256$\,px with a square and a circle that, on hitting a wall, either bounce normally or reverse their velocity with probability one half; initial conditions with a single future or more than ten are excluded. We train on $16{,}384$ clips (one eighth of FlowWM's $131$k) from $5{,}324$ initial conditions with $3.08$ futures on average, and evaluate on $512$ clips from $163$ initial conditions. Because an initial condition is seen with several futures, contexts repeat. Features are FlowWM's DINOv3 definition.

\paragraph{A2: bounce with enumerable directions.}
A ball moves at constant speed in a $64\times64$ frame; at the frame in which it touches a wall, its next direction is drawn uniformly from $K$ equiangular directions spanning $180^\circ$ about the wall normal, and the choice becomes visible one frame later, so about $28\%$ of transitions are stochastic and the rest are deterministic. There are no actions (the action input is zero). \emph{A2-div} draws a new appearance for every episode: background colour, ball colour, ball radius ($3$--$6$\,px), speed ($10$--$18$\,px per frame) and two distractor rectangles of random size and colour. The nearest training context of a test context is $1.9\times$ farther than the typical nearest neighbour within the training set, so contexts never repeat and every training context carries an identifying key. \emph{A2-fix} keeps the appearance constant ($0.35\times$). Each setting has $40$k training frames; coverage is measured on $400$ held-out wall contexts with 16 samples each. An asymmetric variant with $K{=}2$ and probabilities $(0.8, 0.2)$ serves Section~\ref{sec:flatten}.

\paragraph{Push: a branching environment with actions.}
The same renderer with a 2-D push action: each step displaces the ball by $10$\,px along the commanded direction, deflected by $\pm20^\circ$ in the open region (two modes) and by $-60^\circ$, $0^\circ$ or $+60^\circ$ on an ice patch (three modes), with equal probabilities; a pit next to the ice is absorbing. Training uses $40$k frames of random actions. The branch-risk planning task places start and goal just below the pit so that the straight path crosses the ice along the pit's edge.

\paragraph{B: stochastic Push-T.}
We replay the action sequences of the public DINO-WM Push-T data \citep{zhou2025dinowm} in the original simulator, which reproduces the recorded states exactly, and hide one variable per episode: the T block's centre of gravity, drawn from four positions at radius $20$\,px around the default plus $3$\,px of jitter. Under Push-T's quasi-static physics neither mass nor friction changes outcomes, but the centre of gravity changes how the block rotates under the same push, so one pushing context has up to four outcome modes. Five simulator frames form one model step with a 10-dimensional action, as in DINO-WM. Training uses $2{,}400$ episodes ($58{,}398$ frames); the multi-future test set has $128$ starting states taken just before first contact, each with $64$ futures ($4$ modes $\times$ $16$ jitters) of 16 steps; a deterministic control uses the default centre of gravity. A readout decodes the agent position, block position and angle (T-vertex error $2.0$\,px on real features, frame $512$\,px).

\paragraph{BAIR robot pushing.}
$64\times64$ frames with 4-dimensional actions; we train on $4{,}096$ sequences and evaluate on $256$ held-out sequences with 16 samples each, against $1{,}024$ real frames for distribution metrics.

\subsection{Metrics}
\label{sec:metrics}
\emph{Coverage} and \emph{precision} (A2, push): 16 samples per held-out stochastic context are decoded by a position readout (test error $0.13$\,px on real features) and assigned to the nearest enumerated branch within a tolerance; coverage is the fraction of branches hit, precision the fraction of samples within tolerance of some branch. An ideal sampler with 16 draws reaches coverage $1 - (1-1/K)^{16}$, which we report as the ceiling. \emph{Memorization index}: the same quantity on 200 training contexts, restricted to the branch that was realized in the data; chance is $1/K$, and \emph{excess memorization} is $(\text{index} - 1/K)/(1 - 1/K)$. \emph{Mode share}: for the asymmetric variant, the share of samples assigned to the frequent branch. \emph{kNN precision/recall}: the manifold estimates of \citet{kynkaanniemi2019improved} with $k{=}3$ on PCA-reduced tokens. \emph{Deterministic-transition error}: readout error on transitions that are not stochastic. \emph{FlowWM precision/recall error} (A1, B): for each starting state, generate twice as many rollouts as there are true futures; precision error is the mean over rollouts of the distance to the nearest true future, recall error the mean over true futures of the distance to the nearest rollout, and F1 their harmonic mean; distances are L1 over the two shape centres on A1 ($256$\,px frame) and L2 over the eight T vertices on B ($512$\,px). \emph{Manifold distance}: the distance of a predicted token map to the nearest real token map of the test set in a PCA-64 projection, divided by the median real-to-real nearest-neighbour distance, so $1$ means on the data. \emph{Best-of-$n$} and \emph{spread}: readout (or token) error of the sample nearest the truth, and mean pairwise distance of the samples. On BAIR we use token RMS errors, the manifold distance, the spread-to-error ratio, and the Fr\'echet distance and kNN precision/recall of mean-pooled PCA-32 tokens. \emph{Latency}: milliseconds per predicted step at batch size 1. \emph{Planning}: success rate, pit rate and wall-clock time per MPC step.

\subsection{Protocol and statistics}
\label{sec:protocol}
All models are trained with AdamW (learning rate $2\times10^{-4}$, $\beta_2 = 0.9$), batch 128, an EMA of the weights, and $12$k steps unless stated; the budget grid, the asymmetric variant and the per-step-matching comparisons use $8$k steps, iMF $36$k and Shortcut $24$k (Section~\ref{sec:baselines}). Parameter counts are matched within $\pm4\%$ inside each benchmark. Main comparisons use three seeds, the budget grid and the energy-score follow-ups two, ablations one; all means are reported with the standard deviation over seeds. Per-context uncertainty uses a bootstrap over the 400 evaluation contexts, and the Push-T comparisons a hierarchical bootstrap that resamples seeds and then starting states (Appendix~\ref{app:stats}). The experiments were run on two RTX A5000 GPUs, about $300$ GPU-hours in total including the pilot. The predictions for the $K{=}2,16$ grid, for the energy-score follow-ups and for the fed-back-state ablation were written down before the corresponding runs; we reproduce them in the appendix and count the misses.

\section{Results I: Does a Small Candidate Budget Suffice?}
\label{sec:results1}

\begin{figure}[t]
\centering
\includegraphics[width=\linewidth]{fig_scatter_pilot.png}
\caption{Qualitative pilot on the bounce environment with $K{=}4$ branches and fixed appearance (pilot recipe, $m{=}4$). Each column is one context whose ball has just hit a wall: grey, the three context frames; green circles, the four possible next positions; red, 16 samples decoded by the position readout. Regression (top) returns one point near the mean, flow matching with 20 or 2 Euler steps scatters between the modes, and \nom{} places its samples on the modes.}
\label{fig:scatter}
\end{figure}

We first ask whether a one-step \nom{} world model with a small budget covers enumerated futures at all, and how it compares with multi-step and one-step alternatives at matched parameters. Figure~\ref{fig:scatter} shows the phenomenon qualitatively; Tables~\ref{tab:a2}--\ref{tab:b} quantify it on the three benchmarks.

\subsection{Non-repeating contexts (A2-div)}
\label{sec:a2}
Table~\ref{tab:a2} is the main comparison on the bounce environment with $K{=}4$ branches and randomized appearance, where every test context is new (its nearest training context is $1.9\times$ farther than the typical nearest neighbour within the training set). \imlewm{} with $m{=}16$ covers $0.93$ of the enumerated modes against an ideal-sampler bound of $0.99$, with precision $0.82$: its samples lie on the branches. FM-50 reaches $0.79$ coverage but only $0.51$ precision, because its samples scatter between the branches (Figure~\ref{fig:scatter}); FM-2, iMF and Shortcut are below it on both. The energy-score generator has the \emph{lowest} coverage of all stochastic methods ($0.36$ at $m{=}8$, $0.28$ at $m{=}16$) and, like regression, has memorized the training contexts (index $0.87$--$0.97$ against chance $0.25$). \imlewm{} with $m{=}4$ sits in between ($0.46$, index $0.75$). Section~\ref{sec:results2} is about this column. The price of $m{=}16$ is a larger error on transitions that are not stochastic ($1.55$ vs.\ $1.14$\,px for $m{=}4$, $1.09$ for regression) and $1.7\times$ the training time of $m{=}4$; inference cost is independent of $m$, $3.3$\,ms per step against $107$\,ms for FM-50.

\begin{table}[t]
\caption{Bounce with $K{=}4$ branches and randomized appearance (A2-div), 8M-parameter models, $12$k steps, three seeds (mean $\pm$ std; the energy score with $m{=}16$ uses two seeds). Coverage and precision are computed from 16 samples per held-out context decoded to the ball position; the ideal 16-sample sampler reaches coverage $0.99$. kNN recall is the feature-space recall of \citet{kynkaanniemi2019improved}. Det.\ error is the readout error on deterministic transitions. The memorization index is the realized-mode share on training contexts (chance $0.25$). Latency is per predicted step at batch size 1 on a shared RTX A5000.}
\label{tab:a2}
\begin{center}
\small
\resizebox{\linewidth}{!}{%
\begin{tabular}{lcccccccc}
\toprule
Method & NFE & Coverage $\uparrow$ & Precision $\uparrow$ & kNN recall $\uparrow$ & Det.\ err.\ (px) $\downarrow$ & Memo.\ index & ms/step & Train (h) \\
\midrule
Regression & 1 & \ms{0.15}{0.01} & \ms{0.61}{0.02} & 0.00 & \ms{1.09}{0.03} & \ms{0.99}{0.01} & 4.4 & 0.3 \\
FM-50 (FlowWM recipe) & 50 & \ms{0.79}{0.01} & \ms{0.51}{0.00} & \ms{0.91}{0.00} & \ms{2.24}{0.02} & \ms{0.29}{0.01} & 107 & 0.4 \\
FM-2 (same model) & 2 & \ms{0.60}{0.00} & \ms{0.46}{0.01} & \ms{0.56}{0.01} & \ms{1.72}{0.01} & \ms{0.29}{0.01} & 6.5 & 0.4 \\
iMF ($3\times$ steps) & 1 & \ms{0.56}{0.01} & \ms{0.23}{0.02} & \ms{0.92}{0.03} & \ms{14.5}{0.7} & \ms{0.25}{0.01} & 3.4 & 3.2 \\
Shortcut ($2\times$ steps) & 1 & \ms{0.66}{0.00} & \ms{0.42}{0.01} & \ms{0.77}{0.01} & \ms{3.60}{0.07} & \ms{0.27}{0.01} & 4.3 & 0.7 \\
Energy score, $m{=}8$ & 1 & \ms{0.36}{0.02} & \ms{0.61}{0.02} & \ms{0.77}{0.06} & \ms{1.06}{0.04} & \ms{0.87}{0.06} & 3.1 & 1.2 \\
Energy score, $m{=}16$ & 1 & \ms{0.28}{0.01} & \ms{0.61}{0.01} & \ms{0.60}{0.02} & \ms{1.06}{0.03} & \ms{0.97}{0.01} & 3.3 & 2.2 \\
\imlewm{}, $m{=}4$ & 1 & \ms{0.46}{0.01} & \ms{0.68}{0.01} & \ms{0.65}{0.01} & \ms{1.14}{0.07} & \ms{0.75}{0.01} & 3.7 & 0.6 \\
\imlewm{}, $m{=}16$ & 1 & \best{\ms{0.93}{0.01}} & \best{\ms{0.82}{0.01}} & \best{\ms{0.94}{0.01}} & \ms{1.55}{0.05} & \ms{0.29}{0.00} & 3.3 & 1.1 \\
\bottomrule
\end{tabular}}
\end{center}
\end{table}

\subsection{Bouncing Shapes under FlowWM's protocol and readout (A1)}
\label{sec:a1}
Bouncing Shapes \citep{porcher2026flowwm} is the one public benchmark with enumerable futures. We use its generator unchanged, its DINOv3 feature definition and its precision/recall error, but one eighth of its training clips and much smaller models, so we report comparisons within our pipeline and quote the paper's numbers only as a reference. Table~\ref{tab:a1} covers two prediction protocols, autoregressive one-step rollouts (our setting) and the joint prediction of 16 future frames from 16 context frames (FlowWM's setting), and two readouts, our soft-argmax probe and FlowWM's heatmap readout, which we ported. In every one of the four combinations \imlewm{} has the lowest F1 error, and the gaps are at least ten times the seed standard deviation: $6.6$ vs.\ $12.3$\,px for FM-50 in the autoregressive protocol, $9.0$ vs.\ $12.6$\,px with FlowWM's own readout in its own protocol. Because contexts repeat in this benchmark (one initial condition is seen with several futures), there is no memorization and $m{=}4$ is enough; Section~\ref{sec:results2} returns to this difference. The readout matters: the deterministic joint model moves from $5.1$ to $22.5$\,px when the readout changes, which explains most of the gap between our deterministic baseline and the $17.1$\,px reported by \citet{porcher2026flowwm}, while autoregressive models are almost unaffected because they produce a clean feature map at every step (Section~\ref{sec:manifold}). Training four times longer improves every method and preserves the ranking (one seed: \imlewm{} $m{=}4$ $4.1$, FM-2 $7.1$, FM-50 $9.0$\,px).

\begin{table}[t]
\caption{Bouncing Shapes F1 error in px (lower is better; the frame is $256$\,px), three seeds. Autoregressive: 3-frame context, 16 one-step predictions, 8M models. Joint: 16 frames predicted at once from 16 context frames, 27--29M models, 12k steps at batch 64. The last-frame-copy baseline scores $21.7$ in every column. The row from \citet{porcher2026flowwm} uses eight times our training data and a different model size and is a reference only. ``--'' marks combinations that were not run.}
\label{tab:a1}
\begin{center}
\small
\begin{tabular}{lccccc}
\toprule
 & & \multicolumn{2}{c}{Autoregressive} & \multicolumn{2}{c}{Joint 16$\to$16} \\
\cmidrule(lr){3-4}\cmidrule(lr){5-6}
Method & NFE & probe & heatmap & probe & heatmap \\
\midrule
Regression & 1 & \ms{8.09}{0.15} & \ms{8.41}{0.15} & \ms{5.14}{0.01} & \ms{22.48}{0.36} \\
FM-50 & 50 & \ms{12.28}{0.13} & \ms{14.11}{0.20} & \ms{3.86}{0.02} & \ms{12.59}{0.14} \\
FM-2 & 2 & \ms{9.95}{0.19} & \ms{12.49}{0.38} & \ms{3.83}{0.02} & \ms{14.80}{0.17} \\
iMF & 1 & \ms{16.83}{0.24} & -- & -- & -- \\
Shortcut & 1 & \ms{12.07}{0.29} & -- & -- & -- \\
Energy score, $m{=}8$ (2 seeds) & 1 & \ms{8.94}{0.20} & \ms{33.77}{1.34} & -- & -- \\
\imlewm{}, $m{=}4$ & 1 & \best{\ms{6.61}{0.18}} & \best{\ms{6.83}{0.16}} & \best{\ms{3.30}{0.02}} & \ms{11.73}{0.31} \\
\imlewm{}, $m{=}16$ & 1 & \ms{7.21}{0.10} & \ms{7.64}{0.18} & \ms{3.32}{0.07} & \best{\ms{9.04}{0.06}} \\
\midrule
FlowWM paper: FlowWM / deterministic & 50 / 1 & -- & -- & -- & 4.31 / 17.1 \\
\bottomrule
\end{tabular}
\end{center}
\end{table}

\subsection{Stochastic Push-T multi-future test (B)}
\label{sec:b}
On the Push-T variant with a hidden center of gravity, the 21M \imlewm{} with $m{=}16$ has the lowest recall error at both horizons and a precision error $6.6$\,px below FM-50 at $h{=}16$ (Table~\ref{tab:b}). We quantify uncertainty with a hierarchical bootstrap that resamples seeds and then starting states (Appendix~\ref{app:stats}). The precision advantage over FM-50, FM-2 and Shortcut is significant under that procedure ($h{=}16$: $-6.6$ $[-8.5, -4.7]$, $-5.6$ $[-7.7, -3.2]$ and $-13.2$ $[-16.3, -10.2]$\,px). The recall advantage is consistent across the 128 starting states but, with three seeds, its interval includes zero against the flow baselines ($-4.2$ $[-8.8, +1.1]$ vs.\ FM-50) and excludes it only against Shortcut ($-8.2$ $[-11.3, -4.0]$). The evidence on this benchmark is therefore an advantage in precision and a non-inferiority in recall. Two facts limit what any model can show here: all models mispredict the first contact by about $12$\,px (the true block displacement is $16$\,px), and this error does not shrink with three times more training, so differences between objectives are partly masked by a shared dynamics error; and \imlewm{} with $m{=}4$ shows the collapse pattern of Section~\ref{sec:results2}, the best precision with the worst recall among stochastic models and a rollout spread of $9$\,px against $22$\,px for $m{=}16$.

\begin{table}[t]
\caption{Stochastic Push-T multi-future test set (128 pre-contact starting states $\times$ 64 futures $\times$ 16 steps): precision and recall error over the eight T vertices in px (frame $512$\,px), 21M models, three seeds (iMF two). GT spread is the mean pairwise distance among the 64 true futures. Mode coverage is not reported: at these rollout errors no model resolves the four centre-of-gravity modes.}
\label{tab:b}
\begin{center}
\small
\begin{tabular}{lccccc}
\toprule
 & & \multicolumn{2}{c}{$h{=}8$} & \multicolumn{2}{c}{$h{=}16$} \\
\cmidrule(lr){3-4}\cmidrule(lr){5-6}
Method & NFE & precision & recall & precision & recall \\
\midrule
GT spread (reference) & & \multicolumn{2}{c}{30.8} & \multicolumn{2}{c}{48.3} \\
Regression & 1 & \ms{24.8}{1.1} & \ms{39.6}{3.4} & \ms{40.2}{3.0} & \ms{66.4}{8.2} \\
FM-50 & 50 & \ms{27.1}{1.6} & \ms{24.3}{1.8} & \ms{40.7}{1.1} & \ms{42.0}{3.6} \\
FM-2 & 2 & \ms{24.2}{1.7} & \ms{23.3}{1.7} & \ms{39.7}{1.4} & \ms{41.5}{3.8} \\
iMF & 1 & \ms{37.2}{2.6} & \ms{33.5}{3.2} & \ms{55.1}{3.7} & \ms{52.4}{3.4} \\
Shortcut & 1 & \ms{28.0}{1.0} & \ms{24.2}{0.4} & \ms{47.3}{1.9} & \ms{45.9}{1.1} \\
\imlewm{}, $m{=}4$ & 1 & \best{\ms{18.7}{0.4}} & \ms{29.2}{3.4} & \best{\ms{33.8}{3.2}} & \ms{49.1}{3.5} \\
\imlewm{}, $m{=}16$ & 1 & \ms{21.8}{1.8} & \best{\ms{22.1}{2.3}} & \ms{34.1}{0.7} & \best{\ms{37.8}{2.3}} \\
\bottomrule
\end{tabular}
\end{center}
\end{table}

\subsection{Baseline fairness and latency}
\label{sec:fairness}
Four checks guard the comparison. (i) iMF is not under-trained: a 72k-step run on A2-div saturates after 24k steps at coverage $0.58$ and precision $0.26$. (ii) Flow matching is not under-trained: with four times the steps on A1 the ranking is unchanged (Section~\ref{sec:a1}); its recipe is FlowWM's, which on A2-div raised the coverage of our initial flow baseline from $0.58$ to $0.79$. (iii) Shortcut, a stronger one-step flow generator than iMF on every benchmark, is trained to the same wall-clock budget as \imlewm{} with $m{=}16$. (iv) All models within a benchmark are within $\pm4\%$ in parameters. The remaining caveat is scale: our joint FM-50 reaches $12.6$\,px under FlowWM's readout where the paper reports $4.31$ with eight times the data and a larger model, so the claim is confined to matched small-scale compute. Latency at batch size 1 is $3$--$4$\,ms per step for every one-step model and $107$\,ms for FM-50; in the planner of Section~\ref{sec:planning}, where each CEM iteration imagines $64 \times m_p$ futures, this becomes a $50\times$ difference.

\section{Results II: Outcome Memorization and the Candidate Budget}
\label{sec:results2}

\begin{figure}[t]
\centering
\includegraphics[width=\linewidth]{fig_memo_curves.pdf}
\caption{Memorization during training on A2-div ($K{=}4$; mean over seeds, band = min--max). (a) Share of samples on the realized branch of training contexts; chance is $0.25$ and held-out contexts stay at chance for every method. (b) Coverage on held-out contexts. Regression memorizes completely; \nom{} with $m{=}4$ and both energy-score variants memorize progressively and lose coverage as they do; \nom{} with $m{=}16$ and FM-50 do not memorize at this budget.}
\label{fig:curves}
\end{figure}

\subsection{The phenomenon}
\label{sec:phenomenon}
Figure~\ref{fig:curves} tracks the memorization index during training. Regression reaches $0.99$: it learns which of the four branches every training context took and predicts that branch. \imlewm{} with $m{=}4$ climbs to $0.75$ while its held-out coverage falls from $0.77$ at 4k steps to $0.46$ at 12k; on held-out contexts its index stays at chance, so this is not a failure to generalize but the memorization of the realized outcome of each training context, keyed on its appearance. Three further observations fix the interpretation. Training the $m{=}4$ model twice as long raises the index from $0.66$ to $0.92$ and lowers coverage from $0.52$ to $0.38$, and a $2.7\times$ larger model does the same ($0.86$, coverage $0.43$; both with the pilot recipe); the collapse is not under-training. On the same environment with \emph{fixed} appearance (A2-fix), where the identifying key is gone, the same $m{=}4$ recipe does not memorize ($0.26$) and covers $0.97$ of the modes (Table~\ref{tab:es}). And \imlewm{} with $m{=}16$ stays at $0.29$ even when trained $1.5\times$ longer ($0.31$; Appendix~\ref{app:ablation}). The budget, not the recipe, decides.

\subsection{Excess memorization collapses onto $I(m,K)$}
\label{sec:collapse}

\begin{figure}[t]
\centering
\includegraphics[width=\linewidth]{fig_memo_collapse.pdf}
\caption{Excess memorization $(\text{index}-1/K)/(1-1/K)$ on A2-div at 8k steps (two seeds) for $K\in\{2,4,8,16\}$ and $m\in\{2,\dots,64\}$. (a) Against the codebook distortion $I(m,K)$ of Equation~\ref{eq:I}, computed exactly for the bounce geometry. The line and band are the linear fit on the $K{=}4,8$ points ($\pm2$ RMSE) from which the $K{=}2,16$ predictions were recorded before those runs; vertical bars are the recorded intervals. (b) The same points against the miss probability $(1-1/K)^m$, the quantity used by \citet{gladstone2026xm} and \citet{bennett2026freedom}: the points no longer collapse.}
\label{fig:collapse}
\end{figure}

Proposition~\ref{prop:memo} says the incentive to memorize is the random codebook's miss distortion $I(m)$. We trained \imlewm{} on A2-div for $K\in\{4,8\}$ and $m\in\{4,8,16,32\}$ (two seeds, 8k steps), fitted excess memorization linearly in $I(m,K)$ on those eight points (excess $= 0.015 + 1.06\,I$, RMSE $0.032$), and recorded predictions with $\pm2$ RMSE intervals for nine runs at $K\in\{2,16\}$ before training them (Appendix~\ref{app:pred}). Seven of the nine fell inside their interval (Figure~\ref{fig:collapse}a, Table~\ref{tab:grid}). The two misses are at the ends of the range and in the direction we had noted as risks: at $K{=}2$, $m{=}2$ ($I = 0.50$, outside the fitted range) memorization is stronger than the linear extrapolation ($0.96$ vs.\ $0.77$), and at $K{=}16$, $m{=}4$ it is weaker ($0.22$ vs.\ $0.31$), consistent with a capacity effect: storing one of 16 outcomes per context costs $\log_2 16$ bits.

The alternative summary, that the budget should scale with the number of modes, is rejected. At equal $m{=}8$ the excess memorization is $0.075$, $0.083$ and $0.087$ for $K = 4, 8, 16$, as the near-equal $I(8,K)$ on this geometry predicts; at equal $m/K = 1$ it ranges from $0.91$ ($K{=}2$) to $0.06$ ($K{=}16$). A one-variable linear fit on all points except the extrapolated one has RMSE $0.031$ in $I(m,K)$ against $0.068$ in $(1-1/K)^m$ and $0.065$ in $e^{-m/K}$ (Figure~\ref{fig:collapse}b). The linear relation holds for $I \lesssim 0.25$; above that it saturates, and the capacity effect is a second-order factor. The practical rule is that the budget must make $I(m)$ small relative to the mode spacing on the outcome manifold: on this geometry $m = 16$--$32$ brings the excess below $0.05$ for every $K$, and the number of modes enters only through the geometry.

\begin{table}[t]
\caption{Budget grid on A2-div at 8k steps, two seeds. Memorization index (chance $1/K$), coverage, precision and deterministic-transition error (px). Coverage is capped by the ideal 16-sample sampler ($1.00$, $0.99$, $0.88$, $0.64$ for $K = 2, 4, 8, 16$). $K{=}16$ is under-trained at 8k steps: its kNN recall rises from $0.12$ to $0.85$ with $m$ while its coverage stays near $0.28$.}
\label{tab:grid}
\begin{center}
\small
\setlength{\tabcolsep}{4pt}
\begin{tabular}{lccccccccccccc}
\toprule
 & \multicolumn{4}{c}{$K{=}2$} & \multicolumn{4}{c}{$K{=}4$} & \multicolumn{4}{c}{$K{=}8$} \\
\cmidrule(lr){2-5}\cmidrule(lr){6-9}\cmidrule(lr){10-13}
$m$ & 2 & 4 & 8 & 16 & 4 & 8 & 16 & 32 & 4 & 8 & 16 & 32 \\
\midrule
Memo.\ index & 0.96 & 0.57 & 0.53 & 0.52 & 0.50 & 0.31 & 0.27 & 0.26 & 0.31 & 0.20 & 0.15 & 0.14 \\
Coverage & 0.66 & 0.92 & 0.97 & 0.99 & 0.65 & 0.89 & 0.92 & 0.90 & 0.41 & 0.58 & 0.64 & 0.62 \\
Precision & 0.78 & 0.85 & 0.85 & 0.68 & 0.69 & 0.76 & 0.77 & 0.72 & 0.54 & 0.62 & 0.64 & 0.60 \\
Det.\ err.\ (px) & 0.99 & 1.47 & 2.64 & 4.01 & 1.23 & 1.70 & 2.91 & 3.46 & 1.20 & 1.46 & 1.64 & 2.78 \\
\midrule
 & \multicolumn{5}{c}{$K{=}16$} \\
\cmidrule(lr){2-6}
$m$ & 4 & 8 & 16 & 32 & 64 \\
\cmidrule(lr){1-6}
Memo.\ index & 0.22 & 0.14 & 0.12 & 0.11 & 0.08 \\
Coverage & 0.17 & 0.25 & 0.28 & 0.27 & 0.25 \\
Precision & 0.32 & 0.38 & 0.41 & 0.38 & 0.33 \\
Det.\ err.\ (px) & 1.13 & 1.40 & 1.76 & 2.92 & 5.71 \\
\bottomrule
\end{tabular}
\end{center}
\end{table}

\begin{figure}[t]
\centering
\includegraphics[width=\linewidth]{fig_m_knob.pdf}
\caption{What the budget $m$ controls (A2-div, 8k steps, two seeds; shade = $K$). (a)--(c) Coverage and precision peak at $m{=}16$ for $K{=}4,8$ while the memorization index approaches chance. (d)--(e) The costs: the error on deterministic transitions and the training time grow with $m$; inference cost does not. (f) With unequal mode probabilities ($K{=}2$, $p = 0.79$), the share of the frequent mode among \nom{} samples falls with $m$ along the optimum of \citet{bennett2026freedom}; FM-50 is calibrated, and the energy score over-represents the frequent mode (Section~\ref{sec:flatten}).}
\label{fig:knob}
\end{figure}

\subsection{Countermeasures and the cost of $m$}
\label{sec:counter}
Raising $m$ is not the only way to remove the incentive, but it is the cleanest. Table~\ref{tab:counter} compares three alternatives at $m{=}4$ on A2-div ($12$k steps, two seeds). Rejecting candidates that are already within $\epsilon$ of the target, the RS-IMLE prior of \citet{vashist2024rsimle} adapted per transition, changes nothing. Gaussian noise on the context tokens with $\sigma = 1.0$ removes the identifying key and the memorization ($0.27$), but it also blurs the information the prediction needs: precision drops to $0.61$ and the deterministic error rises to $1.92$\,px, and coverage reaches only $0.74$. Noise $\sigma = 0.3$ on top of $m{=}8$ is indistinguishable from clean $m{=}8$. Figure~\ref{fig:knob} shows the costs of $m$ itself: the error on deterministic transitions grows from $1.2$\,px at $m{=}4$ to $2.8$--$3.5$\,px at $m{=}32$ at 8k steps (it shrinks with longer training; $1.55$\,px at 12k for $m{=}16$), and training time grows about threefold from $m{=}4$ to $m{=}32$ because the matching pass scales with $m$ even with late fusion. Coverage and precision saturate at $m{=}16$ on this environment and decline slightly at $m{=}32$. The recipe is otherwise insensitive: replacing the Geman--McClure loss by $\ell_2$, or the 256-dimensional noise by 64 dimensions, changes coverage and the index by at most $0.01$ (Appendix~\ref{app:ablation}).

\begin{table}[t]
\caption{Countermeasures against outcome memorization on A2-div ($K{=}4$, 12k steps, two seeds; the $m{=}4$ and $m{=}16$ rows are the three-seed values of Table~\ref{tab:a2}). The memorization index is given for training / held-out contexts.}
\label{tab:counter}
\begin{center}
\small
\resizebox{\linewidth}{!}{%
\begin{tabular}{lcccccc}
\toprule
Variant & Coverage & Precision & kNN recall & Memo.\ (train / held-out) & Det.\ err.\ (px) & Train (h) \\
\midrule
$m{=}4$ & 0.46 & 0.68 & 0.65 & 0.75 / 0.25 & 1.14 & 0.5 \\
$m{=}4$ + RS-IMLE rejection ($\epsilon{=}0.5$) & 0.46 & 0.69 & 0.65 & 0.74 / 0.25 & 1.17 & 0.5 \\
$m{=}4$ + context noise $\sigma{=}1.0$ & 0.74 & 0.61 & 0.86 & 0.27 / 0.24 & 1.92 & 0.5 \\
$m{=}8$ + context noise $\sigma{=}0.3$ & 0.88 & 0.77 & 0.92 & 0.32 / 0.22 & 1.65 & 0.7 \\
$m{=}16$ & \best{0.93} & \best{0.82} & \best{0.94} & 0.29 / 0.25 & 1.55 & 1.1 \\
\bottomrule
\end{tabular}}
\end{center}
\end{table}

\subsection{A strictly proper score memorizes regardless of $m$}
\label{sec:es}
Proposition~\ref{prop:es} predicts that an energy-score generator, whose memorization incentive does not depend on $m$, memorizes in non-repeating contexts at every budget. We trained the \imlewm{} generator unchanged with the energy score of \citet{makabe2026esp} (Table~\ref{tab:es}). On A2-div it memorizes (index $0.87$ at $m{=}8$, $0.97$ at $m{=}16$; chance $0.25$), its held-out coverage peaks at $0.50$ around 6k steps and then falls to $0.36$ as the index rises (Figure~\ref{fig:curves}), and a larger $m$ makes it memorize \emph{faster}: at 8k steps the index is $0.84$ and $0.87$ for the two $m{=}16$ seeds against $0.51$--$0.72$ for the three $m{=}8$ seeds, as expected if $m$ only lowers the variance of the gradient towards the same minimizer. Replacing the whole-map distance by a per-token energy score, a common choice for continuous tokens, memorizes as well ($0.97$, coverage $0.22$). The ordering of the memorization index across objectives follows the relative incentive of Section~\ref{sec:theory_es}: regression $0.99$, energy score $0.87$--$0.97$, \nom{} with $m = 2, 4, 16$ at $0.96$ ($K{=}2$), $0.75$ and $0.29$. With unequal mode probabilities the score has no calibration advantage either: its frequent-mode share is $0.89$ against a truth of $0.79$, with coverage $0.58$ (Section~\ref{sec:flatten}).

\begin{table}[t]
\caption{Energy score against \nom{} on three bounce settings (12k steps). A2-div: randomized appearance, contexts never repeat. A2-fix: fixed appearance, no identifying key (same recipe for all rows). A1: Bouncing Shapes, contexts repeat; F1 error under the probe and under the FlowWM heatmap readout trained with feature noise $0.05$ and $0.5$. Two seeds unless noted. The one-step manifold distance is the distance of a predicted token map to its nearest real one, in units of the median real-to-real distance.}
\label{tab:es}
\begin{center}
\small
\setlength{\tabcolsep}{4pt}
\resizebox{\linewidth}{!}{%
\begin{tabular}{llcccccc}
\toprule
Setting & Method & Coverage & Precision & kNN recall & Memo.\ index & Det.\ err.\ (px) & Manifold $h{=}1$ \\
\midrule
A2-div & Energy score, $m{=}8$ (3 seeds) & 0.36 & 0.61 & 0.77 & 0.87 & 1.06 & 1.19 \\
($K{=}4$) & Energy score, $m{=}16$ & 0.28 & 0.61 & 0.60 & 0.97 & 1.06 & -- \\
 & Per-token energy score, $m{=}8$ & 0.22 & 0.55 & 0.44 & 0.97 & 1.33 & 1.12 \\
 & \imlewm{}, $m{=}4$ (3 seeds) & 0.46 & 0.68 & 0.65 & 0.75 & 1.14 & 0.84 \\
 & \imlewm{}, $m{=}16$ (3 seeds) & \best{0.93} & \best{0.82} & \best{0.94} & 0.29 & 1.55 & 0.86 \\
\midrule
A2-fix & Energy score, $m{=}8$ & 0.95 & 0.81 & \best{1.00} & 0.27 & 0.20 & $1.59\times$ \\
($K{=}4$) & \imlewm{}, $m{=}4$ & 0.97 & 0.91 & 0.68 & 0.26 & 0.15 & $0.94\times$ \\
 & \imlewm{}, $m{=}16$ & \best{0.99} & \best{0.95} & 0.85 & 0.25 & 0.91 & $1\times$ \\
\midrule
 & & \multicolumn{3}{c}{F1 error (px): probe / heatmap / noisy heatmap} & & & Manifold $h{=}16$ \\
\cmidrule(lr){3-5}
A1 & Energy score, $m{=}8$ & \multicolumn{3}{c}{8.94 / 33.8 / 11.7} & -- & -- & 3.05 \\
 & Regression (3 seeds) & \multicolumn{3}{c}{8.09 / 8.41 / 8.53} & -- & -- & 2.04 \\
 & FM-50 (3 seeds) & \multicolumn{3}{c}{12.28 / 14.11 / 18.38} & -- & -- & 2.95 \\
 & \imlewm{}, $m{=}16$ (3 seeds) & \multicolumn{3}{c}{\best{7.21} / \best{7.64} / \best{7.92}} & -- & -- & 2.08 \\
\bottomrule
\end{tabular}}
\end{center}
\end{table}

\paragraph{Separating the mechanism.}
Memorization needs an incentive, a key that identifies the context, and enough revisits. A toy version of the two-branch problem of \citet{makabe2026esp} ($y = \pm\sqrt{u} + 0.08\varepsilon$, equiprobable branches) with a 16-dimensional nuisance ``appearance'' vector appended to the context reproduces their result when every training context is fresh (energy score matches $0.94$ of samples to a branch centre, \nom{} $0.92$; they report $0.962$ and $0.919$). With a fixed set of 1{,}024 training contexts the energy score memorizes completely ($0.999$) and covers both branches on no new context, while \nom{} with $m = 8$ or $16$ stays at $0.51$--$0.52$ and covers $0.99$--$1.00$ (Figure~\ref{fig:toy}, Appendix~\ref{app:es}). Removing the appearance vector or using $16{,}384$ or more contexts (fewer than $234$ revisits each for this network) removes the memorization of the energy score as well. The two conditions transfer to the real environment: on A2-fix we predicted before the runs that the energy score would not memorize ($\le 0.40$), that it would still trail \imlewm{} with $m{=}16$ in coverage and precision, and that the per-token variant would memorize on A2-div; all five recorded predictions held (Appendix~\ref{app:es}). On A2-fix the coverage gap to \imlewm{} shrinks from $0.57$ to $0.04$ but the precision gap of $0.14$ remains, the energy score's best-of-16 token error is $2.1\times$ larger and its one-step manifold distance $1.6\times$ larger.

\paragraph{The second cause: spread off the manifold.}
In Bouncing Shapes, where contexts repeat and nothing can be memorized, the energy score still trails \imlewm{} and regression (F1 $8.9$ vs.\ $7.2$ and $8.1$\,px). Its error is noise-like: under FlowWM's clean heatmap readout it collapses to $33.8$\,px and under a readout trained with feature noise $0.5$ it recovers to $11.7$, whereas \imlewm{} and regression move by less than $0.5$\,px and flow matching gets worse. Its token-space spread on A2-div is \emph{larger} than that of \imlewm{} with $m{=}16$ (spread-to-error ratio $0.80$ vs.\ $0.64$) while its readout spread is less than half ($1.7$ vs.\ $3.9$\,px) and its rollouts leave the manifold like flow matching does (Table~\ref{tab:es}). The diversity is spent in directions that readouts do not see. This is the behaviour the direction-insensitivity argument of Section~\ref{sec:theory_es} suggests; we present it as an interpretation consistent with the predictions, not as a mechanism we have isolated. The conclusion is conditional, not general: when contexts are fresh or carry no identifying feature the energy score is competitive, as in the toy and on A2-fix; in non-repeating contexts with identifying appearance, the property that \citet{makabe2026esp} list as a defect of the hard minimum, that it is not strictly proper for finite $m$, is what lets \nom{} escape memorization by raising $m$.

\section{Results III: Where No Candidate Budget Suffices}
\label{sec:results3}

\subsection{Mode-probability flattening}
\label{sec:flatten}
Every benchmark so far had equiprobable branches, for which the \nom{} optimum is exactly calibrated (Section~\ref{sec:theory_flatten}). We therefore built a bounce variant with $K{=}2$ branches of probability $0.8$ and $0.2$ (the realized frequency on the evaluation contexts is $0.787$) and measured the share of the frequent branch among 64 samples per context (Table~\ref{tab:flatten}, Figure~\ref{fig:knob}f). The share falls with $m$ exactly as the stationarity condition of \citet{bennett2026freedom} prescribes: $0.82$, $0.65$, $0.57$ for $m = 2, 4, 16$ against theoretical $0.79$, $0.61$, $0.52$. Every method, including FM-50, measures about $0.04$ above its reference, a common offset of the readout-based assignment; relative to FM-50 the measured flattening is $-0.18$ and $-0.26$ against the theoretical $-0.18$ and $-0.27$. The trade-off is therefore real and quantitative. $m{=}2$ is calibrated but memorizes ($0.93$ on training contexts) and covers only $0.75$; $m{=}16$ covers $0.99$ but reports the branches as nearly equiprobable; FM-50 does both. Recovering $p$ from $q$ by inverting the stationarity condition is ill-conditioned for large $m$, since an error in $q$ is amplified in the log-odds by roughly $(m-1)(1/q_1 + 1/q_2)$, about $60\times$ at $m{=}16$, so a large budget should be read as a coverage device, not a probability estimate. For equiprobable branches the calibration is exact, which is why the branch-risk estimates of Section~\ref{sec:planning} are accurate.

\begin{table}[t]
\caption{Asymmetric bounce ($K{=}2$, $p = 0.787$ on the evaluation contexts, 8k steps, two seeds). Share of the frequent branch among 64 samples per held-out context, its theoretical optimum, coverage of both branches, and the realized-branch share on training contexts (chance $0.5$).}
\label{tab:flatten}
\begin{center}
\small
\resizebox{\linewidth}{!}{%
\begin{tabular}{lccccc}
\toprule
Method & Frequent-mode share & Theory & Share $-$ FM-50 (measured / theory) & Coverage & Memo.\ index \\
\midrule
\imlewm{}, $m{=}2$ & 0.823 & 0.787 & $-0.01$ / $0.00$ & 0.75 & 0.93 \\
\imlewm{}, $m{=}4$ & 0.651 & 0.607 & $-0.18$ / $-0.18$ & 0.96 & 0.68 \\
\imlewm{}, $m{=}16$ & 0.573 & 0.522 & $-0.26$ / $-0.27$ & 0.99 & 0.56 \\
FM-50 & 0.832 & 0.787 & -- & 0.95 & 0.73 \\
FM-2 & 0.915 & -- & $+0.08$ / -- & 0.80 & 0.78 \\
Energy score, $m{=}8$ & 0.886 & 0.787 & $+0.05$ / $0.00$ & 0.58 & 0.91 \\
\bottomrule
\end{tabular}}
\end{center}
\end{table}

\subsection{Under-dispersion on real robot video (BAIR)}
\label{sec:bair}
BAIR robot pushing \citep{ebert2017bair} has no enumerable modes: the uncertainty is a continuous, high-dimensional variation of the arm and the objects. Table~\ref{tab:bair} shows the two faces of \nom{} there. On sample accuracy it is the best model: the best-of-16 token error is lowest at every horizon ($0.527$ at $h{=}16$ against $0.843$ for FM-50) and the rollouts stay on the manifold like regression ($1.30$ vs.\ $1.78$). On distribution metrics it loses clearly: Fr\'echet distance $2.09$ against $1.38$ and kNN recall $0.15$ against $0.84$ for FM-50, which needs its 50 steps for this (FM-2 is at the level of regression). The samples are under-dispersed: their spread-to-error ratio at $h{=}1$ is $0.22$ at $m{=}4$ and $0.45$ at $m{=}16$, the same on training and evaluation contexts, so this is not memorization. Raising the budget to $m{=}64$ moves the ratio only to $0.54$ at $2.4\times$ the training time and leaves the Fr\'echet distance at $2.00$; matching in a 512-dimensional pooled-and-projected space, which removes the selection pressure on fine directions, lowers the ratio to $0.31$--$0.38$ and improves recall marginally. No variant within the one-step \nom{} family closes the gap (prediction P4).

\begin{table}[t]
\caption{BAIR robot pushing, 21M models, two seeds (the $m{=}64$ and coarse-matching rows are one seed). Token RMS error of the best of 16 samples and of their mean, manifold distance and spread at $h{=}16$, Fr\'echet distance and kNN recall of mean-pooled PCA-32 tokens against 1{,}024 real frames.}
\label{tab:bair}
\begin{center}
\small
\setlength{\tabcolsep}{4pt}
\resizebox{\linewidth}{!}{%
\begin{tabular}{lccccccccc}
\toprule
Method & NFE & \multicolumn{3}{c}{best-of-16 $\downarrow$ ($h{=}1/8/16$)} & mean $h{=}16$ $\downarrow$ & manifold $h{=}16$ $\downarrow$ & spread $h{=}16$ & FD $h{=}8/16$ $\downarrow$ & recall $h{=}8/16$ $\uparrow$ \\
\midrule
Regression & 1 & 0.268 & 0.481 & 0.556 & 0.556 & \best{1.27} & 0 & 1.44 / 2.65 & 0.21 / 0.08 \\
FM-50 & 50 & 0.338 & 0.681 & 0.843 & 0.885 & 1.78 & 0.94 & \best{0.53 / 1.38} & \best{0.73 / 0.84} \\
FM-2 & 2 & 0.302 & 0.590 & 0.734 & 0.756 & 1.45 & 0.72 & 1.42 / 2.75 & 0.26 / 0.08 \\
Shortcut & 1 & 0.401 & 0.902 & 1.173 & 1.197 & 1.71 & 1.48 & 0.95 / 1.77 & 0.42 / 0.26 \\
\imlewm{}, $m{=}4$ & 1 & 0.260 & 0.460 & 0.532 & \best{0.559} & 1.28 & 0.18 & 1.35 / 2.35 & 0.24 / 0.12 \\
\imlewm{}, $m{=}16$ & 1 & \best{0.258} & \best{0.456} & \best{0.527} & 0.576 & 1.30 & 0.29 & 1.22 / 2.09 & 0.28 / 0.15 \\
\imlewm{}, $m{=}64$ & 1 & 0.257 & 0.461 & 0.534 & 0.589 & 1.33 & 0.32 & 1.23 / 2.00 & 0.28 / 0.14 \\
\imlewm{}, $m{=}16$, coarse matching & 1 & 0.259 & 0.463 & 0.535 & 0.565 & 1.29 & 0.21 & 1.21 / 2.03 & 0.33 / 0.15 \\
\bottomrule
\end{tabular}}
\end{center}
\end{table}

\begin{figure}[t]
\centering
\includegraphics[width=\linewidth]{fig_bair.pdf}
\caption{The shape of the under-dispersion on BAIR ($h{=}1$, 256 evaluation contexts). (a) Spread-to-error ratio of the samples against $m$. The isotropic model of Proposition~\ref{prop:shrink} with one free parameter fitted at $m{=}16$ ($\deff \approx 27$) reproduces the $m{=}4$ and $m{=}64$ points within $0.05$; with $d$ set to the independently measured participation ratio of the regression residuals ($\approx 1{,}030$) it predicts a spread five times too small. (b) Participation ratio of the 64 samples of each context (median and quartiles): the generator spreads along one to three directions per context. (c) Fraction of the true residual that lies in a context's own top-16 sample directions versus in another context's directions: the directions are context-specific and aligned with the truth, but capture only $9$--$21\%$ of it.}
\label{fig:bair}
\end{figure}

\paragraph{The shape of the contraction.}
Proposition~\ref{prop:shrink} reproduces the growth of the spread with $m$ if one fits its single parameter: with $\deff \approx 27$ chosen to match $m{=}16$, the predicted ratios at $m = 4$ and $64$ are $0.27$ and $0.58$ against measured $0.22$ and $0.54$ (Figure~\ref{fig:bair}a). We tried to remove that free parameter by measuring the effective dimension independently, as the participation ratio of the residuals of the regression model in token space, and recorded the prediction that it would fall within a factor of two of $27$. It did not: the pooled residual is extremely high-dimensional (participation ratio $\approx 1{,}030$; its largest direction carries $0.6\%$ of the variance) and would predict a spread five times smaller than observed. The isotropic model is therefore a one-parameter description, not a parameter-free prediction. An exploratory analysis designed after this failure explains the discrepancy: per context, the generator's samples occupy only $1.2$, $1.8$ and $2.5$ effective directions at $m = 4, 16, 64$ (Figure~\ref{fig:bair}b), these directions differ across contexts, and they are aligned with the true residual of their own context ten times better than with that of another context, while capturing only $9$--$21\%$ of it (Figure~\ref{fig:bair}c). The contraction is not an isotropic shrinkage but a collapse onto a few context-specific directions whose number grows slowly with $m$, the shape of reverse water-filling in rate--distortion theory, where a rate of $\ln m$ nats is spent on the largest-variance components first and components below the water level are dropped. We could not test this quantitatively because each context has a single true future. Together with Section~\ref{sec:flatten} this delimits the method: a candidate budget fixes coverage of a few discrete modes, it does not reproduce a continuous high-dimensional conditional distribution, and it distorts mode probabilities in between.

\section{Results IV: Rollouts and Planning}
\label{sec:results4}

\subsection{Rollouts stay on the token manifold}
\label{sec:manifold}

\begin{figure}[t]
\centering
\includegraphics[width=\linewidth]{fig_manifold.pdf}
\caption{Manifold distance of autoregressive rollouts against the horizon (16 rollouts from real contexts; three seeds; log scale). The distance of a predicted token map to its nearest real token map is divided by the median real-to-real nearest-neighbour distance, so $1$ means on the data. \nom{} rollouts stay as close to the data as deterministic regression on all three benchmarks; flow-matching rollouts drift by the same amount with 50 and with 2 integration steps; iMF and the energy score drift further.}
\label{fig:manifold}
\end{figure}

A one-step world model is used recursively, so what matters beyond one step is whether its outputs remain inputs the model has seen. Figure~\ref{fig:manifold} measures this as the normalized distance of each predicted token map to the nearest real one. On A2 the distance of \imlewm{} with $m{=}16$ grows from $0.86$ at $h{=}1$ to $1.68$ at $h{=}32$, within $0.1$ of regression ($1.80$), while FM-50 reaches $6.16$ and FM-2 $6.36$, so the drift is not an integration error; iMF reaches $8.04$. The same ordering holds on Bouncing Shapes ($2.08$ vs.\ $2.95$--$3.57$ at $h{=}16$) and Push-T ($3.28$ vs.\ $4.37$--$5.81$). Staying on the manifold is why the autoregressive models in Table~\ref{tab:a1} are nearly indifferent to the readout while the joint-prediction models, which emit 16 frames at once with less precise per-frame features, lose a factor of $2.7$--$4.4$ under the heatmap readout. It also explains the one benchmark where FM has the better best-of-16 error: on Push-T the flow rollouts spread more ($36$ vs.\ $22$\,px at $h{=}16$) while drifting, so one of them lands closer to the truth ($67$ vs.\ $79$\,px); on A1 and A2 \imlewm{} is best on both. A control rules out readout fragility as the cause of the flow baselines' deficit: under a heatmap readout trained with feature noise $0.5$, regression and \imlewm{} improve ($22.5 \to 9.7$ and $9.0 \to 5.3$\,px in the joint protocol) and flow matching gets worse ($12.6 \to 15.0$); its errors are structural, not noise.

\subsection{Per-step matching on self-rollouts and what the fed-back state does}
\label{sec:seq}
Per-step matching (Section~\ref{sec:seqmatch}) trains on the model's own rollout with $H{=}4$. On stochastic Push-T it reduces the $h{=}16$ recall error from $43.7 \pm 0.9$ to $29.3 \pm 2.8$\,px (per-seed change $-33 \pm 8\%$, all three seeds improve by 10--19\,px; hierarchical bootstrap $-14.3$ $[-19.1, -9.9]$\,px) and the precision error by $41\%$, at 8k steps for both; it beats the one-step-matched model trained $1.5\times$ longer (Table~\ref{tab:seq}). On the deterministic Push-T control, where the only train--test gap is exposure bias, the best-of-16 error at $h{=}16$ falls by $19\%$ (one seed), so most of the gain is the classic remedy working. On A2 only the manifold distance improves ($-24\%$ at $h{=}32$) while one-step coverage and precision drop by $0.03$--$0.04$; on A1 the best-of-16 rollout error improves by $24\%$ and recall by $9\%$ while precision worsens by $6\%$, leaving F1 unchanged. The gain is therefore largest where rollout error compounds through contact dynamics. The cost is $3$--$4\times$ the training time; inference is unchanged.

\begin{table}[t]
\caption{Per-step matching on self-rollouts ($H{=}4$) against one-step matching at equal training (8k steps). Left: stochastic Push-T, 21M, three seeds, precision/recall error in px with the paired per-seed change. Right: the fed-back-state ablation on Bouncing Shapes (8M, $m{=}16$, two seeds): manifold distance, best-of-16 readout error and spread of 16 rollouts at $h{=}16$ (true futures spread $23.6$\,px), and F1 error. Predictions S1--S3 recorded before the ablation all held (Appendix~\ref{app:seq}).}
\label{tab:seq}
\begin{center}
\scriptsize
\setlength{\tabcolsep}{4pt}
\begin{tabular}{lcccc}
\toprule
Stochastic Push-T & $h{=}8$ prec.\ / recall & $h{=}16$ prec.\ / recall & Train (h) \\
\midrule
One-step matching & 25.7 / 25.4 & 41.6 / 43.7 & 1.3 \\
Per-step, feed matched & \best{18.8} / \best{18.4} & \best{24.6} / \best{29.3} & 5.3 \\
Paired change & $-27\%$ / $-28\%$ & $-41\%$ / $-33\%$ & $\times 3.9$ \\
One-step matching, 12k steps & 21.8 / 22.1 & 34.1 / 37.8 & 2.2 \\
\bottomrule
\end{tabular}
\hfill
\begin{tabular}{lcccc}
\toprule
Bouncing Shapes, fed-back state & Manifold & Best-of-16 & Spread & F1 \\
\midrule
One-step matching & 2.17 & 23.2 & 17.0 & 8.74 \\
Per-step, feed matched & 1.92 & \best{17.7} & 26.3 & \best{8.51} \\
Per-step, feed random candidate & \best{1.81} & 22.9 & 12.4 & 8.77 \\
Per-step, feed true frame & 2.14 & 23.6 & 16.6 & 8.62 \\
\bottomrule
\end{tabular}
\end{center}
\end{table}

\begin{figure}[t]
\centering
\includegraphics[width=\linewidth]{fig_seq_ablation.pdf}
\caption{The fed-back-state ablation on Bouncing Shapes (two seeds). (a) Any self-generated context, matched or random, keeps the rollout on the manifold; the true frame does not. (b) Only the candidate nearest the truth improves the best-of-16 error. (c) The matched variant is the only one whose 16 rollouts spread like the true futures at $h{=}16$; the random variant jitters at every step but does not branch, and spreads to half the truth.}
\label{fig:seqabl}
\end{figure}

Selecting the candidate nearest the ground truth is the device of closest-among-top-$K$ rollouts \citep{zhang2025catk}; we asked which part of the gain it accounts for by changing only the fed-back state (Table~\ref{tab:seq}, Figure~\ref{fig:seqabl}). Feeding back the \emph{true} frame, which performs the same number of matched updates on real contexts, leaves everything at the one-step-matching level, so the gain is not a batch-size or update-count effect. Feeding back a \emph{random} candidate keeps the rollout on the manifold at least as well as the matched one ($1.81$ vs.\ $1.92$) but does not improve the best-of-16 error ($22.9$ vs.\ $23.2$) and halves the spread of the rollouts relative to the true futures ($0.52\times$): when the fed-back candidate sits on a different branch than the truth, the true next frame becomes an unreachable target, the one-step distribution widens into per-step jitter, and the jitter does not accumulate into branches. Only the matched candidate keeps the context on the branch of the truth, so the branch continues ($1.11\times$ the true spread at $h{=}16$, best-of-16 $17.7$). Two effects are thus separable: the manifold benefit comes from training on any self-generated context, the coverage benefit from feeding back the candidate nearest the truth, which in \nom{} training is the argmin of the loss itself. The ablation was run on Bouncing Shapes only.

\subsection{Planning: equal success, 50$\times$ faster, accurate branch risk}
\label{sec:planning}
A one-step model lets a planner imagine many futures per candidate action sequence at low cost, and a model that covers the futures should in principle avoid the risky ones. We tested both with CEM-MPC on two tasks (Table~\ref{tab:plan}). On stochastic Push-T (100 episodes, three seeds; success if the block ends within $20$\,px and $20^\circ$ of the goal) \imlewm{}, regression and FM-2 are indistinguishable at $32$--$35\%$, FM-50 is at $30\%$ on ten episodes, and the planning time per step is $4.0$\,s for \imlewm{}, about half of FM-2 and $1/50$ of FM-50. The contact error of Section~\ref{sec:b} is the likely reason no model exploits the hidden centre of gravity.

\begin{table}[t]
\caption{Planning with CEM-MPC. Left: stochastic Push-T, 100 episodes $\times$ 3 seeds (FM-50: 10 episodes, 1 seed), planning time per MPC step. Right: branch-risk task (push along the pit edge), 100 episodes $\times$ 3 seeds, with the model-estimated pit probability after one push on the ice (true $0.33$; three equiprobable deflections, one of which falls into the pit) and after five straight pushes from 200 start states (true $0.43$).}
\label{tab:plan}
\begin{center}
\scriptsize
\setlength{\tabcolsep}{4pt}
\begin{tabular}{lcc}
\toprule
Stochastic Push-T & Success & s / step \\
\midrule
Regression & \ms{31.7}{3.3}\% & 4.3 \\
FM-2 & \ms{34.7}{3.1}\% & 8.5 \\
FM-50 & 30\% & 200 \\
\imlewm{}, $m{=}16$ & \ms{31.7}{2.1}\% & \best{4.0} \\
\bottomrule
\end{tabular}
\hfill
\begin{tabular}{lcccc}
\toprule
Branch-risk task & Success & Pit & One-push pit prob.\ (true 0.33) & Five-push pit prob.\ (true 0.43) \\
\midrule
Oracle (true simulator) & 76\% & 7\% & 0.33 & 0.43 \\
Regression & \ms{56.3}{3.4}\% & \ms{24.3}{2.5}\% & \ms{0.61}{0.10} & \ms{0.26}{0.19} \\
FM-2 & \best{\ms{67.0}{1.4}\%} & \best{\ms{18.7}{0.9}\%} & \ms{0.14}{0.00} & \ms{0.38}{0.06} \\
FM-50 (estimates only) & -- & -- & \ms{0.31}{0.02} & 0.44 \\
\imlewm{}, $m{=}16$ & \ms{55.3}{2.6}\% & \ms{31.0}{2.2}\% & \best{\ms{0.36}{0.03}} & \best{\ms{0.44}{0.04}} \\
\bottomrule
\end{tabular}
\end{center}
\end{table}

The branch-risk task isolates the question. The straight path from start to goal crosses an ice patch along the edge of an absorbing pit; each push on the ice is deflected by $-60^\circ$, $0^\circ$ or $+60^\circ$ with equal probability, and the upward deflection falls in. A world model that represents the three branches can estimate the risk of the straight path. \imlewm{} does: its one-push pit probability is $0.36 \pm 0.03$ against the true $0.33$, and its five-push estimate $0.44$ against $0.43$, the most accurate of all models (per-start absolute error $0.078$). FM-2 under-estimates the one-push risk by $2.4\times$ ($0.14$) because two Euler steps do not resolve the branches; the same model with 50 steps gets it right ($0.31$) at $24\times$ the planning time. Regression cannot represent it ($0.61$, with a per-start error of $0.46$). Yet the planner built on \imlewm{} succeeds no more often than the one built on regression ($55\%$ vs.\ $56\%$) and falls into the pit \emph{more} often ($31\%$ vs.\ $24\%$); the FM-2 planner, with the worst risk estimates, succeeds most often ($67\%$; the oracle planner reaches $76\%$). Planning performance does not follow the model's probabilistic accuracy. The bottleneck is the planner: a mean-distance cost over a handful of imagined futures lets the cross-entropy method exploit model errors, and an absorbing failure state is not penalized much more than a detour. This decoupling agrees with recent reports that one-step fidelity predicts planning poorly \citep{cheung2026fidelity,terver2025drives} and that one-step world models bring no planning gain on Push-T \citep{lee2026beckmann}. A planner that uses the branch probabilities directly, through a conditional value-at-risk or a failure-probability constraint, is the natural next step; we did not build one. A mean-plus-variance cost is not the answer: with the true simulator as the world model it lowered success from $78\%$ to $50\%$ as the penalty grew (Appendix~\ref{app:negative}), because a certain failure has zero variance.

\section{Discussion and Limitations}
\label{sec:discussion}

\paragraph{When to use a one-step \nom{} world model.}
The results reduce to one criterion. If the outcome uncertainty consists of a few discrete alternatives, as in bouncing directions, a hidden centre of gravity or a branching obstacle, the random codebook can cover them: the distortion $I(m)$ vanishes exponentially once $m$ exceeds the number of modes on the outcome manifold, a budget of $16$--$32$ removes the memorization incentive, and the model covers the modes with one forward pass, stays on the token manifold under rollout and estimates equiprobable branch risks correctly. If instead the uncertainty is continuous and high-dimensional, as in the texture and fine motion of real video, the codebook cannot cover it at any affordable $m$; the samples collapse onto a few context-specific directions and the generator is best read as an accurate point predictor with a little aligned spread. In between, the budget distorts mode probabilities in a known way. A practical rule for choosing $m$ is to estimate $I(m)$ from the outcome diversity of the training data, for instance from a simulator that enumerates branches or from the outcomes of near-identical contexts, and to pick the smallest $m$ at which $I(m)$ is a few percent of the mode spacing. A hybrid that uses \nom{} to select the mode and a small stochastic head or a one- or two-step flow for the residual is the natural way to cover both regimes; we leave it to future work.

\paragraph{Limitations.}
(1) The approach under-disperses on real video and no variant in the one-step family fixed it; the quantitative agreement of the isotropic model on BAIR uses one fitted parameter, and the independent estimate of that parameter failed. (2) On stochastic Push-T the advantage over flow matching is an advantage in precision; the recall advantage does not clear seed variability with three seeds, and a shared $12$\,px contact error masks differences. (3) The benchmarks are synthetic or simulated except BAIR; the models are small, and our flow baselines, although trained with FlowWM's recipe and checked for under-training, are far from the scale of \citet{porcher2026flowwm}, so the comparison is confined to matched small-scale compute. (4) Planning success did not improve on either task, even where the model's risk estimates were accurate. (5) Large budgets flatten mode probabilities; the samples of such a model are a coverage device and need a separate calibration mechanism, which the energy score did not provide in non-repeating contexts. (6) The ablations, the BAIR variants and the deterministic Push-T control are single-seed; the $K{=}2,16$ grid and the energy-score follow-ups are two-seed. (7) Proposition~\ref{prop:es} was derived after we observed the energy score memorizing; only the follow-up runs (A2-fix, per-token score) had predictions recorded in advance, and we examined one $\beta$ and two distance choices, not mixtures of the two objectives. (8) The mathematics of Section~\ref{sec:theory} is classical; the contribution is the application and its tests.

\section{Conclusion}
\label{sec:conclusion}
A one-step world model trained with a \nom{} objective on frozen DINO tokens covers enumerable futures better than multi-step flow matching and other one-step generators at matched parameters, with one network evaluation per step and rollouts that stay on the data manifold. The objective is the distortion of a random codebook of $m$ draws, and that single fact answers the question in the title: the budget must make the codebook's miss distortion small relative to the spacing of the outcome modes, which a few dozen candidates achieve for mode-like uncertainty and no budget achieves for continuous high-dimensional uncertainty. The same accounting shows why, in contexts that never repeat, a strictly proper energy score memorizes the realized outcome while \nom{} can be made not to, why mode probabilities flatten with the budget, and where the method ends. We hope the memorization index and the enumerable benchmarks are useful to others building one-step stochastic world models, whichever objective they choose.


\subsubsection*{Broader Impact Statement}
This work studies the training objective of small world models on synthetic and simulated data and on a public robot-pushing dataset. We do not see a significant risk of harm beyond those of world modeling in general. One point deserves a note: \nom{} models with a large budget over-represent rare modes (Section~\ref{sec:flatten}), so their samples must not be read as calibrated probabilities in safety-relevant use.

\subsubsection*{Reproducibility}
All environments, training and evaluation code, the queue files that define every run, and the per-run outputs used for the figures and tables will be released. Appendix~\ref{app:impl} lists hyperparameters, parameter counts and training budgets; Appendix~\ref{app:bench} specifies the benchmarks; Appendix~\ref{app:stats} describes the bootstrap procedures. The predictions recorded before the $K{=}2,16$ runs and before the energy-score follow-up runs are reproduced verbatim in Appendices~\ref{app:pred} and~\ref{app:es}.

\bibliography{main}
\bibliographystyle{tmlr}

\appendix
\section{Derivations}
\label{app:proofs}

\paragraph{Proposition~\ref{prop:memo}.}
With non-repeating contexts, the empirical loss of a conditional generator decomposes over training contexts. For the memorizing solution every candidate equals the realized mode, so the minimum distance is zero for every context. For the generalizing solution the candidates are $m$ i.i.d.\ draws from $p(\cdot\,|\,c_i)$, and the realized mode is itself a draw from $p$, so the expected minimum distance is $I(m)$ of Equation~\ref{eq:I}. The difference of the two empirical risks is therefore $I(m)$. For the bounds, let the modes be equiprobable and separated by distances in $[\delta, \Delta]$. The probability that none of the $m$ draws equals the realized mode $k$ is $(1 - 1/K)^m$; on that event the nearest draw is at squared distance in $[\delta^2, \Delta^2]$, and otherwise the distance is zero, which gives the stated bounds. If the modes densely fill a curve of length $\ell$ ($K \gg m$), the draws are $m$ uniform points on the curve and the expected squared arc distance from a uniform point to its nearest draw is $\Theta(\ell^2/m^2)$, as for the nearest of $m$ uniform points on an interval, so $I(m)$ decays polynomially and does not depend on $K$. For squared-error regression the single prediction can memorize the realized mode at zero loss, and the generalizing prediction (the conditional mean) has loss $\mathrm{tr}\,\mathrm{Cov}_p$, independent of any budget. A detail used in Section~\ref{sec:collapse}: the realized-mode share $s$ relates to the excess memorization through $s = 1/K + (1 - 1/K)\,\mathrm{excess}$, so a $\pm 2$ RMSE interval on the excess scale maps to $\pm 2(1 - 1/K)\,\mathrm{RMSE}$ on the index scale.

\paragraph{Proposition~\ref{prop:shrink}.}
Write $Y_j - X = (\mu - X) + s\,\varepsilon_j$ with $\varepsilon_j \sim \mathcal{N}(0, I_d)$ and $X - \mu \sim \mathcal{N}(0, \sigma^2 I_d)$. Then $\|Y_j - X\|^2 = \|X - \mu\|^2 + s^2 \|\varepsilon_j\|^2 - 2 s\, \varepsilon_j^\top (X - \mu)$. The first term is common to all $j$; the second concentrates at $s^2 d$ with fluctuations of order $s^2 \sqrt{d}$; the third is, conditionally on $X$, Gaussian with variance $4 s^2 \|X - \mu\|^2 \approx 4 s^2 \sigma^2 d$ and independent across $j$, so the minimum over $j$ of the sum is, to leading order in $d$, $d\sigma^2 + d s^2 - 2 s \sigma \sqrt{d}\, e_m$ with $e_m = \E \max_{j \le m} \xi_j$. Minimizing in $s$ gives $s^\star = \sigma e_m / \sqrt{d}$. The spread-to-error ratio used in Section~\ref{sec:bair}, defined as the mean pairwise distance among samples divided by the mean distance from samples to the truth, is $r = \sqrt{2}\rho/\sqrt{1 + \rho^2}$ with $\rho = s/\sigma$ under this model, which is how $\deff$ is fitted from one measured $r$.

\paragraph{Flattening.}
Let the generator place mass $q_k$ on mode $k$ and let all modes be at mutual distance $\Delta$. The probability that none of $m$ draws lands on the realized mode $k$ is $(1 - q_k)^m$, so $L(q) = \Delta^2 \sum_k p_k (1 - q_k)^m$. Minimizing subject to $\sum_k q_k = 1$ gives the stationarity condition $m\, p_k (1 - q_k)^{m-1} = \lambda$ for every $k$ with $q_k > 0$, which is Theorem~7 of \citet{bennett2026freedom}; for $K = 2$ it reads $\big((1-q)/q\big)^{m-1} = (1-p)/p$. If $p_k = 1/K$ for all $k$ the symmetric point $q_k = 1/K$ satisfies the condition and is the unique minimizer of the convex objective for $m \ge 2$ (Theorem~4 of the same work). The numerical optima quoted in the text (for $p = (0.787, 0.213)$: $0.787$, $0.607$, $0.522$ at $m = 2, 4, 16$) solve this equation.

\paragraph{Proposition~\ref{prop:es}.}
(a) The energy distance between $q$ and $\delta_x$ is $2\E_{Y \sim q}\|Y - x\|^\beta - \E_{Y, Y' \sim q}\|Y - Y'\|^\beta - 0$, so $\IS_\beta(q, x)$ is half of it; the energy distance with $\beta \in (0, 2)$ is a metric on distributions with finite moments \citep{gneiting2007strictly,dawid2007geometry}, hence non-negative and zero only when $q = \delta_x$. The empirical risk over training contexts is a sum of such terms with one $x$ per context, so its unique minimizer is the memorizing solution; the $m$-sample estimator with the unbiased pairwise term has this risk as its expectation for every $m \ge 2$. (b) At $q = p$, $\E_{x \sim p}\,\IS_\beta(p, x) = \E\|Y - x\|^\beta - \frac{1}{2}\E\|Y - Y'\|^\beta = \frac{1}{2}\E\|Y - Y'\|^\beta$ because $x$ and $Y$ are exchangeable. (c) For $\beta = 2$, $\E\|Y - x\|^2 - \frac{1}{2}\E\|Y - Y'\|^2 = \|\E Y - x\|^2$. (d) The per-token score is a sum over tokens of energy scores of the token marginals, and (a)--(b) apply to each term.

\paragraph{Direction insensitivity (interpretation).}
Suppose the generator's mean misses the outcome by an unknown vector of length $a$ and it adds isotropic noise of scale $s$ in $d$ dimensions with $d$ large. Then $\E\|Y - x\|^\beta \approx (a^2 + s^2)^{\beta/2}$ and $\E\|Y - Y'\|^\beta \approx (2 s^2)^{\beta/2}$, and the energy score $(a^2 + s^2)^{\beta/2} - \frac{1}{2}(2s^2)^{\beta/2}$ is minimized at $s = a$ for every $\beta \in (0, 2)$, where it equals $2^{\beta/2 - 1} a^\beta$; this is the $\lambda = 1$ case of Proposition~4.1 of \citet{debortoli2025distributional}, and the gain does not depend on $d$. The same calculation for the \nom{} loss follows Proposition~\ref{prop:shrink} and gives $s \approx a\, e_m/\sqrt{d}$. The approximation requires the bias and the noise to be spread over many dimensions (participation ratio $\gg 1$); the regression residual on BAIR has participation ratio about $1{,}030$, and we did not measure it on the bounce tokens. The argument says nothing about where the isotropic spread points; the claim in Section~\ref{sec:es} that it points off the manifold is an empirical observation.

\section{$I(m,K)$ for the bounce geometry and the recorded predictions}
\label{app:pred}
\label{app:imk}
At a wall frame the $K$ possible next directions are the unit vectors at angles $\theta_k = \big((k + \tfrac12)/K\big)\,180^\circ - 90^\circ$ about the wall normal, $k = 0, \dots, K-1$, with unit speed. For a target mode $k$ let $d^{(k)}_{(1)} \le \dots \le d^{(k)}_{(K)}$ be the sorted squared distances from $\mu_k$ to all modes ($d^{(k)}_{(1)} = 0$). The nearest of $m$ i.i.d.\ uniform draws is the $r$-th closest mode with probability $\big((K - r + 1)/K\big)^m - \big((K - r)/K\big)^m$, so
\begin{equation}
I(m, K) = \frac{1}{K} \sum_{k=1}^{K} \sum_{r=1}^{K} d^{(k)}_{(r)} \Big[ \big(\tfrac{K - r + 1}{K}\big)^m - \big(\tfrac{K - r}{K}\big)^m \Big].
\end{equation}
Table~\ref{tab:predfull} lists the values together with the predictions recorded before the $K \in \{2, 16\}$ runs. The predictions were made by fitting excess memorization linearly in $I$ on the eight $K \in \{4, 8\}$ points (excess $= 0.015 + 1.06\,I$, RMSE $0.032$) and mapping to the index scale; the intervals are $\pm 2$ RMSE. We noted two risks at the time: the $K{=}2$, $m{=}2$ point lies outside the fitted range of $I$, and large $K$ might memorize less than predicted because more bits must be stored per context. Both misses are of these kinds.

\begin{table}[h]
\caption{Recorded predictions for the memorization index at 8k steps and the observed two-seed means. The two misses are marked.}
\label{tab:predfull}
\begin{center}
\small
\resizebox{\linewidth}{!}{%
\begin{tabular}{cccccc}
\toprule
$K$ & $m$ & $I(m,K)$ & Chance & Predicted [interval] & Observed \\
\midrule
2 & 2 & 0.500 & 0.50 & 0.77 [0.74, 0.81] & 0.956 (miss) \\
2 & 4 & 0.125 & 0.50 & 0.57 [0.54, 0.61] & 0.568 \\
2 & 8 & 0.008 & 0.50 & 0.51 [0.50, 0.54] & 0.528 \\
2 & 16 & $\approx 0$ & 0.50 & 0.51 [0.50, 0.54] & 0.522 \\
16 & 4 & 0.234 & 0.0625 & 0.31 [0.25, 0.37] & 0.221 (miss) \\
16 & 8 & 0.074 & 0.0625 & 0.15 [0.09, 0.21] & 0.144 \\
16 & 16 & 0.021 & 0.0625 & 0.10 [0.06, 0.16] & 0.115 \\
16 & 32 & 0.005 & 0.0625 & 0.08 [0.06, 0.14] & 0.105 \\
16 & 64 & 0.0006 & 0.0625 & 0.08 [0.06, 0.14] & 0.081 \\
\midrule
4 & 4 / 8 / 16 / 32 & 0.235 / 0.061 / 0.006 / 0.0001 & 0.25 & (fitted) & 0.498 / 0.306 / 0.272 / 0.263 \\
8 & 4 / 8 / 16 / 32 & 0.236 / 0.076 / 0.019 / 0.002 & 0.125 & (fitted) & 0.305 / 0.198 / 0.154 / 0.144 \\
\bottomrule
\end{tabular}}
\end{center}
\end{table}

\section{Benchmark construction}
\label{app:bench}

\paragraph{A2 appearance randomization.}
Each episode of A2-div draws a 22-dimensional appearance vector: background RGB in $[10, 110]$, ball RGB in $[150, 255]$, radius in $[3, 6]$\,px, speed in $[10, 18]$\,px per frame, and two rectangles with corner in $[0, 48]$, size in $[6, 20]$ and RGB in $[40, 200]$. Frames are rendered at $64\times64$ with an anti-aliased disc, upsampled to $224$\,px and encoded by DINOv2 ViT-S/14. Context novelty is measured as the ratio of the median distance from a test context to its nearest training context to the median nearest-neighbour distance within the training set in a PCA-64 token space: $1.9$ for A2-div and $0.35$ for A2-fix. The evaluation uses 400 held-out wall contexts; a stochastic transition is one whose current frame touches a wall.

\paragraph{Push environment.}
The $64\times64$ arena contains an ice rectangle ($x \in [20, 44]$, $y \in [14, 40]$) and an absorbing pit ($x \in [24, 40]$, $y \in [0, 14]$). A push action is a 2-D vector; with probability $0.1$ it is zero, otherwise its direction is uniform and its magnitude in $[0.3, 1]$. The ball moves $10$\,px along the action rotated by the sampled deflection and respawns after four frames in the pit. The branch-risk task draws starts with $x \in [8, 12]$, $y \in [17, 21]$ and goals with $x \in [52, 56]$ in the same band. The planner uses CEM-MPC with horizon $T{=}5$, 64 candidate sequences, $m_p{=}8$ imagined futures per candidate (one for regression), and a cost equal to the mean distance between the decoded final position and the goal; 100 episodes per model and seed with a shared set of task seeds. The risk diagnostic evaluates 200 start states with the straight five-push action sequence against 256 true-simulator rollouts per start, and a one-push test from 200 ice states just below the pit.

\paragraph{Stochastic Push-T hidden variable.}
Table~\ref{tab:cog} shows the per-episode randomizations we tried before choosing the centre of gravity; mass and friction have no effect under Push-T's quasi-static physics, and damping produces unrealistic sliding. The multi-future test set takes 128 starting states just before the first contact (an earlier version with earlier starts did not diverge within 8 steps and was discarded), and for each the 64 futures share the three context frames. Modes closer than $2$\,px in outcome are merged, leaving $3.9$ distinguishable modes per start on average. The CEM planner for Push-T uses $T{=}5$, 64 candidates, $m_p{=}4$, 3 iterations with 8 elites, re-planning every step for at most 10 steps, with the mean token-space $\ell_2$ distance to the goal as the cost, following \citet{terver2025drives}; goals are within 25 steps of the start, and success requires the block within $20$\,px and $20^\circ$ of the goal.

\begin{table}[h]
\caption{Candidate hidden variables for stochastic Push-T: deviation of the T vertices after 16 steps under the same actions (median / maximum over episodes; the nominal displacement is $124$\,px).}
\label{tab:cog}
\begin{center}
\small
\begin{tabular}{lcl}
\toprule
Randomization & Deviation (px) & Decision \\
\midrule
Block mass $0.5$--$2\times$ & 0.0 / 0.0 & no effect (quasi-static) \\
Contact friction $0.3$--$1.2$ & 0.0 / 0.0 & no effect \\
Damping $0$--$0.1$ & 25 / 346 & unrealistic sliding \\
Centre of gravity, radius $20$\,px, 4 positions + 3\,px jitter & 28 / 138 & adopted \\
Centre of gravity, radius $25$\,px & 56 / 195 & too large \\
\bottomrule
\end{tabular}
\end{center}
\end{table}

\paragraph{BAIR.}
We use the first part of the official archive: $4{,}096$ training sequences and $1{,}024$ held-out sequences (the official test set plus three held-out training files), $64\times64$ frames, 4-dimensional actions, 3 context frames and 16-step rollouts.

\section{Implementation details and ablations}
\label{app:impl}

\paragraph{Training.}
AdamW with learning rate $2\times10^{-4}$, $(\beta_1, \beta_2) = (0.9, 0.9)$, no weight decay, a linear warm-up over 1{,}000 steps followed by a cosine decay, batch 128, gradient clipping at norm 1, an exponential moving average of the weights with decay $\min(0.999, (1 + t)/(10 + t))$ used for all evaluations, 12k steps (8k for the budget grid, the asymmetric variant and the per-step-matching comparisons; 36k for iMF; 24k for Shortcut; 48k for the four-times-longer A1 runs). The noise vector has 256 dimensions and is normalized to unit RMS before the mapping network. Late fusion sets the first five of eight blocks (eight of twelve for the 21M model) to a learned constant style. Candidate generation for matching runs under no-gradient in bfloat16 in a separate autocast region. The Geman--McClure scale is twice the running median (EMA $0.99$) of the matching distance. The energy-score runs use $\beta = 1$, the whole-map RMS distance, the unbiased pairwise estimator, and backpropagate in two chunks of contexts; the per-token variant computes pairwise differences explicitly to be exact near zero. Flow matching uses the time shift $\alpha = 13$ at training and sampling and the Euler sampler. The iMF port follows the official JAX implementation through a PyTorch re-implementation: dual $u$/$v$ heads of width 640 (to match parameters), the JVP direction taken from the model's own $v$ prediction, and the loss $\|u + (t - r)\,\mathrm{sg}(\partial_t u) - v\|^2$ plus a flow-matching loss on $v$; it converges only with the time-embedding frequency scale set to 1 rather than the usual 1000. Our own MeanFlow implementation diverged three times once the average-velocity term was switched on and was abandoned. Shortcut models devote a quarter of each batch to the self-consistency target computed with the EMA model. Parameter counts: 8M setting, $7.97$--$8.24$M across methods ($8.24$M for \imlewm{}); 21M setting, $20.6$M for \imlewm{}. Training times in the tables are wall-clock hours on one RTX A5000; some runs shared a GPU and are slower than the stand-alone figures (8M, batch 128, one step: $m{=}4$ 178\,ms, $m{=}16$ 488\,ms; 21M: 404 and 1112\,ms).

\paragraph{Recipe choices.}
Table~\ref{tab:recipe} gives the pilot comparisons behind the recipe. The wide head does not change \imlewm{} (it is there to match FM); late fusion halves training time at equal coverage and lowers the deterministic error; the FlowWM recipe raised the flow baseline substantially. Multi-stage supervision, which ROMS-IMLE found important for images, does not help on a fixed $16\times16$ token grid (Appendix~\ref{app:negative}).

\begin{table}[h]
\caption{Recipe comparisons on A2-div ($K{=}4$, pilot and preparation runs, one seed; bootstrap 95\% intervals over 400 contexts where given). Training hours are on a shared GPU.}
\label{tab:recipe}
\begin{center}
\small
\resizebox{\linewidth}{!}{%
\begin{tabular}{lcccccc}
\toprule
Variant & Coverage & Precision & kNN recall & Det.\ err.\ (px) & Memo.\ (train / val) & Train (h) \\
\midrule
\imlewm{} $m{=}16$, full noise injection, wide head & 0.934 [0.917, 0.949] & 0.83 & 0.94 & 2.31 & 0.29 / 0.24 & 1.93 \\
\imlewm{} $m{=}16$, late fusion, wide head (adopted) & 0.925 [0.907, 0.942] & 0.83 & 0.95 & 1.56 & 0.29 / 0.27 & 1.05 \\
\imlewm{} $m{=}16$, full injection, narrow head (6.6M) & 0.930 [0.909, 0.948] & 0.84 & 0.94 & 2.02 & 0.30 / 0.25 & 1.71 \\
\imlewm{} $m{=}16$, pilot recipe (three spatial stages) & 0.91 [0.89, 0.93] & 0.81 & 0.94 & 2.67 & 0.30 / 0.22 & 3.2 \\
FM, default recipe, 20 Euler steps & 0.58 [0.55, 0.61] & 0.38 & 0.68 & 4.64 & -- & 0.6 \\
FM, FlowWM recipe, 50 steps (adopted) & 0.79 [0.77, 0.82] & 0.51 & 0.91 & 2.27 & 0.29 / 0.24 & 0.4 \\
iMF, time-embedding scale 1000 & \multicolumn{5}{c}{does not converge} \\
iMF, time-embedding scale 1, 12k steps & 0.55 [0.52, 0.58] & 0.22 & -- & 13.1 & -- & 1.0 \\
\bottomrule
\end{tabular}}
\end{center}
\end{table}

\paragraph{Ablations.}
\label{app:ablation}
On A2-div with $m{=}16$ (one seed): replacing Geman--McClure by $\ell_2$ gives coverage $0.93$, precision $0.83$, index $0.28$, deterministic error $1.81$\,px (vs.\ $1.55$); a 64-dimensional noise vector gives $0.92$, $0.83$, $0.28$, $1.50$; training $1.5\times$ longer (18k steps) gives $0.92$, $0.86$, $0.31$, $1.45$. On stochastic Push-T an 8M \imlewm{} with $m{=}16$ reaches recall errors $24.5$ / $39.0$\,px at $h = 8$ / $16$ against $22.1$ / $37.8$ for the 21M model (three seeds), so model size matters little there. In the pilot with the earlier recipe, Geman--McClure improved precision ($0.92$ vs.\ $0.86$) and token-space coverage ($0.85$ vs.\ $0.66$) over $\ell_2$. The iMF learning curve on A2-div (one seed, checkpoints of a 72k-step run) saturates after 24k steps: coverage $0.56$, $0.57$, $0.58$, $0.58$ and precision $0.24$, $0.25$, $0.26$, $0.26$ at 24k, 42k, 60k and 72k steps.

\section{Readouts}
\label{app:readout}
Position readouts are small networks trained on real tokens: a soft-argmax probe for the ball (test error $0.13$\,px on A2-div, $0.05$\,px with DINOv3 on A1), and for Push-T a probe for agent position, block position and angle (T-vertex error $2.0$\,px). Because a probe might decode generated (out-of-distribution) tokens differently from real ones, the pilot cross-checked it against a nearest-neighbour judge that decodes a token map by looking up the nearest of $12{,}000$ real frames in a PCA-128 space: on the fixed-appearance environment the two judges agreed within $0.01$ for every \imlewm{} run (e.g.\ coverage $0.960$ vs.\ $0.960$ at $K{=}4$) and the nearest-neighbour judge was $0.03$--$0.09$ lower for flow matching ($0.841$ vs.\ $0.814$ with 20 steps), so the probe is, if anything, generous to the baselines; the ranking was identical under both. The FlowWM heatmap readout, ported from the public code, maps the $384$-channel map through a $1\times1$ convolution to 256 channels, GELU, dropout $0.1$ and a $16\times$ transposed convolution to $256\times256$ logits per shape, trained with pixel cross-entropy and feature noise of standard deviation $0.05$ for 20k steps, and reads the argmax; its error on real features is $0.10$\,px (p95 $0.75$). The noise-robust variant is trained with feature noise $0.5$ (floor $1.03$\,px). Table~\ref{tab:readout} gives all readout combinations for A1.

\begin{table}[h]
\caption{Bouncing Shapes F1 error (px) under three readouts, three seeds (energy score two). The last row is the 48k-step run of Section~\ref{sec:a1} (one seed).}
\label{tab:readout}
\begin{center}
\small
\begin{tabular}{lcccccc}
\toprule
 & \multicolumn{3}{c}{Joint 16$\to$16} & \multicolumn{3}{c}{Autoregressive} \\
\cmidrule(lr){2-4}\cmidrule(lr){5-7}
Method & probe & heatmap & heatmap, noise 0.5 & probe & heatmap & heatmap, noise 0.5 \\
\midrule
Regression & 5.14 & 22.48 & 9.67 & 8.09 & 8.41 & 8.53 \\
FM-50 & 3.86 & 12.59 & 15.00 & 12.28 & 14.11 & 18.38 \\
FM-2 & 3.83 & 14.80 & 16.43 & 9.95 & 12.49 & 16.85 \\
Energy score, $m{=}8$ & -- & -- & -- & 8.94 & 33.77 & 11.65 \\
\imlewm{}, $m{=}4$ & 3.30 & 11.73 & 5.92 & 6.61 & 6.83 & 7.04 \\
\imlewm{}, $m{=}16$ & 3.32 & 9.04 & 5.28 & 7.21 & 7.64 & 7.92 \\
\imlewm{}, $m{=}4$, 48k steps & -- & -- & -- & 4.08 & 4.20 & -- \\
\bottomrule
\end{tabular}
\end{center}
\end{table}

\section{Additional results}
\label{app:more}

\subsection{Rollout metrics}
Table~\ref{tab:e3} completes Figure~\ref{fig:manifold} with the best-of-16 readout error and the spread at the last horizon.

\begin{table}[h]
\caption{Rollouts from real contexts (16 samples, three seeds; iMF on Push-T two): manifold distance, best-of-16 readout error (px) and spread (px) at the last horizon ($h{=}32$ on A2, $16$ on A1 and B).}
\label{tab:e3}
\begin{center}
\small
\setlength{\tabcolsep}{4pt}
\begin{tabular}{lccccccccc}
\toprule
 & \multicolumn{3}{c}{A2 bounce, $h{=}32$} & \multicolumn{3}{c}{A1 Bouncing Shapes, $h{=}16$} & \multicolumn{3}{c}{B Push-T, $h{=}16$} \\
\cmidrule(lr){2-4}\cmidrule(lr){5-7}\cmidrule(lr){8-10}
Method & manifold & best-of-16 & spread & manifold & best-of-16 & spread & manifold & best-of-16 & spread \\
\midrule
Regression & 1.80 & 32.2 & 0 & 2.04 & 34.0 & 0 & 3.21 & 102.2 & 0 \\
FM-50 & 6.16 & 9.8 & 25.7 & 2.95 & 28.0 & 25.1 & 4.60 & \best{67.4} & 36.4 \\
FM-2 & 6.36 & 10.6 & 26.8 & 3.03 & 23.4 & 26.1 & 4.37 & 69.2 & 29.7 \\
iMF & 8.04 & 10.1 & 29.0 & 3.57 & 33.4 & -- & 5.81 & 84.1 & 57.0 \\
\imlewm{}, $m{=}4$ & 1.74 & 10.5 & 23.9 & \best{1.98} & \best{19.0} & -- & \best{3.20} & 88.7 & 9.2 \\
\imlewm{}, $m{=}16$ & \best{1.68} & \best{9.3} & 27.2 & 2.08 & 19.4 & 17.9 & 3.28 & 79.1 & 21.8 \\
\bottomrule
\end{tabular}
\end{center}
\end{table}

\subsection{Stochastic Push-T: bootstrap intervals}
Table~\ref{tab:bci} gives the paired differences behind Section~\ref{sec:b} and Section~\ref{sec:seq}.

\begin{table}[h]
\caption{Paired differences on the stochastic Push-T multi-future test set (px; negative favours the first model) with 95\% bootstrap intervals over the 128 starting states after averaging seeds, and hierarchical intervals that also resample the three seeds (10{,}000 resamples).}
\label{tab:bci}
\begin{center}
\small
\begin{tabular}{llcccc}
\toprule
Comparison & Metric & Difference & Start-state CI & Hierarchical CI \\
\midrule
\imlewm{} 16 vs.\ FM-50 & $h{=}8$ recall & $-1.8$ & $[-2.5, -1.1]$ & $[-5.2, +2.3]$ \\
 & $h{=}8$ precision & $-5.2$ & $[-5.7, -4.8]$ & $[-8.0, -2.5]$ \\
 & $h{=}16$ recall & $-4.2$ & $[-5.3, -3.1]$ & $[-8.8, +1.1]$ \\
 & $h{=}16$ precision & $-6.6$ & $[-7.6, -5.7]$ & $[-8.5, -4.7]$ \\
\imlewm{} 16 vs.\ FM-2 & $h{=}16$ recall & $-3.7$ & $[-4.9, -2.5]$ & $[-8.7, +1.9]$ \\
 & $h{=}16$ precision & $-5.6$ & $[-6.6, -4.6]$ & $[-7.7, -3.2]$ \\
\imlewm{} 16 vs.\ Shortcut & $h{=}16$ recall & $-8.2$ & $[-10.0, -6.3]$ & $[-11.3, -4.0]$ \\
 & $h{=}16$ precision & $-13.2$ & $[-15.1, -11.3]$ & $[-16.3, -10.2]$ \\
Per-step vs.\ one-step matching (8k) & $h{=}8$ recall & $-7.1$ & $[-8.2, -6.0]$ & $[-9.2, -4.8]$ \\
 & $h{=}16$ recall & $-14.3$ & $[-16.3, -12.5]$ & $[-19.1, -9.9]$ \\
 & $h{=}16$ precision & $-17.0$ & $[-18.4, -15.6]$ & $[-20.4, -14.5]$ \\
\bottomrule
\end{tabular}
\end{center}
\end{table}

\subsection{Energy-score follow-ups}
\label{app:es}
Figure~\ref{fig:toy} shows the toy experiment of Section~\ref{sec:es}. The task is the two-branch problem of \citet{makabe2026esp} ($u \sim U[0.05, 1]$, $y = \pm\sqrt{u} + 0.08\varepsilon$, equiprobable branches) with a 16-dimensional nuisance vector $a \sim \mathcal{N}(0, I)$ appended to the context, a three-layer MLP generator of width 256 with a 4-dimensional noise input, Adam at $2\times10^{-4}$, batch 256, 15k steps, and the same energy-score estimator as the main runs. With fresh contexts at every step the energy score matches $0.940$ / $0.946$ of its samples to a branch centre within $0.24$ ($m = 8$ / $16$) and \nom{} $0.916$ / $0.924$. With 1{,}024 fixed training contexts (about 3{,}750 revisits each) the energy score memorizes ($0.999$) and covers both branches on $0.2\%$ of new contexts, while \nom{} with $m = 8$ / $16$ is at $0.52$ / $0.51$ and covers $0.99$ / $1.00$. With 4{,}096, 16{,}384 and 65{,}536 contexts the energy score's realized-branch share is $0.73$, $0.51$ and $0.50$; without the nuisance vector (context $= u$ only) it is $0.50$ even with 1{,}024 contexts. The thresholds are those of this small network and do not transfer; the 8M model on A2-div saw each of its $11{,}559$ stochastic training contexts about 38 times and memorized.

\begin{figure}[h]
\centering
\includegraphics[width=0.9\linewidth]{fig_toy.pdf}
\caption{Toy two-branch problem with a nuisance appearance vector (two seeds, $m{=}8$ for the energy score and $16$ for \nom{}). (a) Realized-branch share on training contexts against the number of fixed training contexts. (b) Two-branch coverage on new contexts; ``fresh'' draws a new context at every step.}
\label{fig:toy}
\end{figure}

The follow-up runs on A2-fix and with the per-token score were preceded by five recorded predictions, all of which held (two seeds each): (P1) the energy score with $m{=}8$ does not memorize on A2-fix (index $\le 0.40$; observed $0.268$); (P2) its coverage on A2-fix exceeds its A2-div value but stays below \imlewm{} with $m{=}16$ in both coverage and precision (observed $0.947$ vs.\ $0.987$ and $0.814$ vs.\ $0.954$); (P3) \imlewm{} with $m{=}4$ does not memorize on A2-fix ($\le 0.35$; observed $0.256$); (P4) the per-token energy score memorizes on A2-div (index $\ge 0.7$, coverage $\le 0.5$; observed $0.970$ and $0.224$); (P5) its one-step manifold distance exceeds $1.1$ (observed $1.12$; \imlewm{} $0.86$). On A2-fix the energy score's kNN recall is $1.00$ against $0.85$ for \imlewm{}, its samples spread more widely than the outcomes; its deterministic error is the lowest ($0.20$\,px). On A2-div its learning curves peak in coverage at 6k steps ($0.50$) and decline as the index rises ($0.42 \to 0.87$); the two $m{=}16$ seeds reach an index of $0.84$ and $0.87$ at 8k steps against $0.51$--$0.72$ for the three $m{=}8$ seeds. Token-space diagnostics at $h{=}1$ on A2-div (three seeds, evaluation contexts): spread-to-error ratio $0.80$ (energy score) vs.\ $0.64$ (\imlewm{} 16), $0.38$ (\imlewm{} 4), $1.15$ (FM-50); best-of-16 token error $0.468$ vs.\ $0.271$; mean-prediction error $0.451$ vs.\ $0.396$; readout spread $1.7$ vs.\ $3.9$\,px; one-step manifold distance $1.19$ vs.\ $0.86$.

\subsection{Per-step matching}
\label{app:seq}
Per-seed Push-T recall errors at $h{=}16$ (one-step $\to$ per-step): $43.7 \to 30.3$, $42.6 \to 32.2$, $44.7 \to 25.5$. On the deterministic Push-T control (one seed) the best-of-16 T-vertex error at $h = 4 / 8 / 16$ moves from $35.8 / 60.4 / 72.3$ to $29.2 / 45.6 / 58.9$\,px and the manifold distance at $h{=}16$ from $3.90$ to $3.64$. On A2-div (two seeds, 8k steps) the manifold distance at $h = 8 / 16 / 32$ moves from $1.31 / 1.47 / 1.83$ to $1.19 / 1.26 / 1.40$ while one-step coverage and precision move from $0.92 / 0.77$ to $0.90 / 0.73$ and the deterministic error from $2.91$ to $3.97$\,px; the best-of-16 error is unchanged. On A1 (two seeds) F1 moves from $8.74$ to $8.51$ (precision $9.26 \to 9.77$, recall $8.28 \to 7.55$), the manifold distance at $h{=}16$ from $2.17$ to $1.92$ and the best-of-16 error from $23.2$ to $17.7$\,px. The fed-back-state ablation was preceded by three recorded predictions: (S1) feeding back the true frame yields less than half of the manifold gain (manifold at $h{=}16$ $\ge 2.05$; observed $2.14$); (S2) feeding back a random candidate yields more than half ($\le 2.05$; observed $1.81$); (S3) the random candidate's best-of-16 error exceeds the matched one's by at least $1$\,px ($\ge 18.7$; observed $22.9$). The spread of the 16 rollouts at $h = 1 / 4 / 8 / 16$ is $0.20 / 1.90 / 6.81 / 23.6$\,px for the true futures, $0.70 / 3.35 / 7.42 / 17.0$ for one-step matching, $1.16 / 5.42 / 12.2 / 26.3$ for the matched variant, $2.41 / 4.19 / 6.38 / 12.4$ for the random variant and $0.61 / 3.07 / 7.08 / 16.6$ for the true-frame variant. Training times: one-step $0.68$\,h, per-step variants $2.7$\,h (8M, 8k steps).

\subsection{Planning details}
On stochastic Push-T the final block position and angle errors are $46.6$\,px / $42.0^\circ$ (regression), $40.6$ / $35.2$ (FM-2), $45.2$ / $20.0$ (FM-50, 10 episodes) and $49.5$ / $44.9$ (\imlewm{}). The 1-step contact error (block displacement after a contact transition, true $15.9$\,px) is $11.7$ / $11.2$\,px for regression at 12k / 36k steps, $12.4$ / $12.6$ for FM-2 / FM-50 and $12.2$ / $12.3$ for \imlewm{} with $m = 4$ / $16$ (seed 0). On the branch-risk task the per-start absolute error of the five-push pit probability is $0.078$ (\imlewm{}), $0.086$ (FM-50), $0.096$ (FM-2) and $0.462$ (regression), and the timeout rates are $13.7\%$, $14.3\%$, $19.3\%$ and $17\%$ (oracle). Wall-clock for 100 episodes: oracle $100$\,s, regression $523$\,s ($m_p{=}1$), \imlewm{} $3{,}633$\,s, FM-2 $9{,}067$\,s.

\section{Negative results}
\label{app:negative}
The pilot study tested four additions proposed in our original design; none survived.

\paragraph{Trajectory-level matching.}
Matching $m{=}4$ candidate rollouts of length $H{=}8$ as whole trajectories (with a curriculum $H = 1 \to 4 \to 8$) on the push environment was worse than one-step matching at every horizon: best-of-16 position error $1.31$ vs.\ $0.35$\,px at $h{=}1$, $3.89$ vs.\ $2.19$ at $h{=}8$, $4.67$ vs.\ $3.03$ at $h{=}16$ (paired bootstrap intervals exclude zero), one-step coverage $0.71$ vs.\ $0.98$, and $3.3\times$ the training time; only the error of the sample mean improved ($9.18$ vs.\ $9.94$ at $h{=}16$), the signature of averaging. With two to three branches per step an 8-step trajectory has hundreds of modes, far beyond $m{=}4$, so the matched candidate is the one close to all of them on average. Per-step matching (Section~\ref{sec:seqmatch}) replaced it.

\paragraph{Multi-stage supervision.}
On the fixed $16\times16$ token grid, supervising intermediate heads with spatially pooled targets ($4\times4$, $8\times8$, $16\times16$) or with PCA-truncated targets (32, 128, 384 components) gave the same coverage as a single head within the bootstrap interval ($0.960$ and $0.959$ vs.\ $0.967$ at $K{=}4$) and the same precision ($0.916$, $0.922$ vs.\ $0.926$). The final recipe uses a single head.

\paragraph{Variance-penalized planning.}
A CEM cost of mean plus $\lambda$ times the spread of the imagined goal distances lowered the success rate of the branch-risk task even when the true simulator served as the world model: $78\%$, $77\%$, $70\%$ and $50\%$ at $\lambda = 0, 0.5, 1, 2$, with the pit rate rising to $35\%$. A certain failure has zero variance, so the penalty prefers it to an uncertain success. With \imlewm{} as the model (two seeds, 60 episodes each) the pooled success was $62.5\%$ at $\lambda = 0$ and $61.7\%$ at $\lambda = 0.5$.

\paragraph{Counterfactual action contrast.}
We had anticipated that the noise $z$ might absorb the effect of the action. It does not: on the push environment the counterfactual precision (samples under a random other action staying away from the true outcome) and the action sensitivity (output change per action change relative to the truth) are $0.95$ and $0.96$ with the action on its own FiLM path, $0.92$ and $0.97$ with the action merged into the noise path, and $0.93$ and $0.975$ with an added margin term, against $0.83$ and $0.94$ for regression. There was no shortcut to repair.

\paragraph{Low-dimensional matching and larger budgets on BAIR.}
See Table~\ref{tab:bair}: $m{=}64$ raises the spread-to-error ratio from $0.44$--$0.51$ to $0.53$--$0.55$ at $2.4\times$ the training time, and matching in a 512-dimensional pooled space lowers it to $0.31$--$0.38$; neither approaches FM-50 on distribution metrics.

\section{Statistical methodology}
\label{app:stats}
Per-context intervals are percentile bootstrap intervals (2{,}000 resamples) over the 400 evaluation contexts of a run, used in the pilot and recipe tables. Seed statistics are means and standard deviations over the listed seeds. For the Push-T multi-future test set, each model's per-start errors were stored; the start-state interval averages the seeds first and resamples the 128 starts, and the hierarchical interval resamples the seeds with replacement and then the starts, 10{,}000 times each. Per-step and one-step matching runs with the same seed are paired. All comparisons report the direction and the magnitude of the difference; where an interval includes zero we say so in the text.


\end{document}
